\documentclass[11pt]{article}

\usepackage[a4paper,margin=24mm,headheight=15pt,footskip=14mm]{geometry}
\usepackage{fontspec}
\defaultfontfeatures{Ligatures=TeX,Scale=MatchLowercase}
\setmainfont{Latin Modern Roman}
\setsansfont{Latin Modern Sans}
\setmonofont{Latin Modern Mono}
\usepackage{microtype}
\usepackage{amsmath,amssymb,mathtools}
\usepackage{booktabs,longtable,tabularx,array,multirow}
\usepackage{enumitem}
\usepackage{xcolor}
\usepackage{hyperref}
\usepackage{xurl}
\usepackage{bookmark}
\usepackage{listings}
\usepackage{fancyhdr}
\usepackage{titlesec}
\usepackage{caption}
\usepackage{float}
\usepackage{ragged2e}
\usepackage{needspace}
\usepackage{parskip}

\definecolor{Ink}{HTML}{202428}
\definecolor{Muted}{HTML}{5F6B73}
\definecolor{Accent}{HTML}{315F72}
\definecolor{AccentLight}{HTML}{EAF1F4}
\definecolor{Rule}{HTML}{BCC9CF}
\definecolor{CodeBg}{HTML}{F4F6F7}
\definecolor{CodeFrame}{HTML}{CDD5D9}
\definecolor{LinkBlue}{HTML}{1B5E82}
\definecolor{WarnBg}{HTML}{FFF7E8}
\definecolor{GoodBg}{HTML}{EDF6EE}

\hypersetup{
  colorlinks=true,
  linkcolor=Accent,
  citecolor=Accent,
  urlcolor=LinkBlue,
  pdfauthor={OpenAI},
  pdftitle={Future Directions for Djex and Leant},
  pdfsubject={Co-evolution and automation of term synthesis for complex types},
  pdfkeywords={Djex, Leant, Lean 4, Djinn, Exference, term synthesis, rank-N types}
}
\urlstyle{same}

\pagestyle{fancy}
\fancyhf{}
\fancyhead[L]{\small\sffamily Djex + Leant: future directions}
\fancyhead[R]{\small\sffamily Review snapshot: 2026-08-10}
\fancyfoot[C]{\small\sffamily\thepage}
\renewcommand{\headrulewidth}{0.35pt}
\renewcommand{\footrulewidth}{0pt}

\titleformat{\section}
  {\Large\bfseries\sffamily\color{Ink}}
  {\thesection}{0.7em}{}
\titleformat{\subsection}
  {\large\bfseries\sffamily\color{Ink}}
  {\thesubsection}{0.7em}{}
\titleformat{\subsubsection}
  {\normalsize\bfseries\sffamily\color{Ink}}
  {\thesubsubsection}{0.7em}{}
\titlespacing*{\section}{0pt}{2.5ex plus .8ex minus .3ex}{1.1ex}
\titlespacing*{\subsection}{0pt}{2.0ex plus .6ex minus .2ex}{0.8ex}
\titlespacing*{\subsubsection}{0pt}{1.6ex plus .5ex minus .2ex}{0.6ex}

\setlist[itemize]{leftmargin=1.5em,itemsep=0.28em,topsep=0.35em}
\setlist[enumerate]{leftmargin=1.65em,itemsep=0.32em,topsep=0.4em}
\setlist[description]{leftmargin=2.0em,labelindent=0pt,itemsep=0.35em,style=nextline}

\newcolumntype{Y}{>{\RaggedRight\arraybackslash}X}
\newcolumntype{L}[1]{>{\RaggedRight\arraybackslash}p{#1}}
\newcolumntype{C}[1]{>{\Centering\arraybackslash}p{#1}}

\lstdefinestyle{code}{
  basicstyle=\small\ttfamily,
  backgroundcolor=\color{CodeBg},
  frame=single,
  rulecolor=\color{CodeFrame},
  framerule=0.4pt,
  xleftmargin=0.7em,
  xrightmargin=0.4em,
  framexleftmargin=0.35em,
  framexrightmargin=0.2em,
  aboveskip=0.8em,
  belowskip=0.8em,
  breaklines=true,
  breakatwhitespace=false,
  columns=fullflexible,
  keepspaces=true,
  showstringspaces=false,
  tabsize=2,
  keywordstyle=\bfseries\color{Accent},
  commentstyle=\itshape\color{Muted},
  stringstyle=\color{LinkBlue}
}
\lstset{style=code}

\newcommand{\code}[1]{\texttt{#1}}
\newcommand{\Djex}{\textsf{Djex}}
\newcommand{\Leant}{\textsf{Leant}}
\newcommand{\Lean}{Lean~4}
\newcommand{\Djinn}{Djinn}
\newcommand{\Exference}{Exference}
\newcommand{\RepoDjex}{\href{https://github.com/VladimirReshetnikov/Djex}{\code{VladimirReshetnikov/Djex}}}
\newcommand{\RepoLeant}{\href{https://github.com/VladimirReshetnikov/Leant}{\code{VladimirReshetnikov/Leant}}}
\newcommand{\djcommit}[2]{\href{https://github.com/VladimirReshetnikov/Djex/commit/#1}{\code{#2}}}
\newcommand{\lecommit}[2]{\href{https://github.com/VladimirReshetnikov/Leant/commit/#1}{\code{#2}}}
\newcommand{\headsha}[1]{\code{#1}}
\newcommand{\term}[1]{\emph{#1}}
\newcommand{\yes}{\textbf{yes}}
\newcommand{\no}{\textbf{no}}
\newcommand{\risk}[1]{\textsf{#1}}

\newenvironment{finding}[1]
{\par\medskip\noindent\colorbox{AccentLight}{\parbox{\dimexpr\linewidth-2\fboxsep\relax}{\sffamily\bfseries #1}}\par\smallskip}
{\par\medskip}

\newenvironment{recommendation}[1]
{\par\medskip\noindent\colorbox{GoodBg}{\parbox{\dimexpr\linewidth-2\fboxsep\relax}{\sffamily\bfseries #1}}\par\smallskip}
{\par\medskip}

\newenvironment{caution}[1]
{\par\medskip\noindent\colorbox{WarnBg}{\parbox{\dimexpr\linewidth-2\fboxsep\relax}{\sffamily\bfseries #1}}\par\smallskip}
{\par\medskip}

\setlength{\emergencystretch}{3em}
\sloppy

\begin{document}
\hypersetup{pageanchor=false}

\begin{titlepage}
  \centering
  \vspace*{1.5cm}
  {\sffamily\bfseries\fontsize{25}{31}\selectfont
   From Bounded Rank-N Plumbing to\\[0.2em]
   Typed, Feedback-Directed Synthesis\par}
  \vspace{0.9cm}
  {\sffamily\Large Future directions for \Djex{} and \Leant{}\par}
  \vspace{0.45cm}
  {\sffamily\large A review of their coordinated git history through 10 August 2026\par}
  \vspace{1.5cm}
  \rule{0.78\textwidth}{0.55pt}
  \vspace{1.2cm}

  \begin{minipage}{0.83\textwidth}
  \large
  \textbf{Repositories reviewed}\par
  \vspace{0.25cm}
  \RepoDjex\\
  \RepoLeant

  \vspace{0.75cm}
  \textbf{Review snapshot}\par
  \vspace{0.25cm}
  \Djex{} main: \headsha{f50b5eac86726c7530d95e48c27101b7c0fa25c9}\\
  \Leant{} main: \headsha{0160237e4bdf6e342a55f5a86ce5d5292728024f}\\
  \Leant{} submodule \code{lib/Djex}: the same \Djex{} head
  \end{minipage}

  \vfill
  {\sffamily\small Technical report prepared 10 August 2026\par}
\end{titlepage}

\hypersetup{pageanchor=true}
\pagenumbering{roman}

\begin{abstract}
\noindent
The recent histories of \Djex{} and \Leant{} show an unusually productive
co-development loop.  \Leant{} exposes a concrete Lean-side synthesis failure;
that failure identifies information erased at the language-neutral boundary;
\Djex{} then generalizes the missing invariant or search rule; and \Leant{}
advances its pinned submodule, preserves the source semantics that remain
Lean-specific, and closes the loop with kernel-checked transcripts.  In roughly
the last two days alone, this loop widened bounded instantiation to six leading
binders, carried higher-kinded class arguments, preserved contextual rank-N
provider assignments, enforced assignment fidelity through nested and sibling
datatype fields, specialized provider results from exact term evidence, and
finally retained the structural identities of \code{Prod}/\code{Sum} through
higher-kinded substitutions.

The success of that loop also reveals its next scaling limit.  More and more
semantic information is transported beside an intentionally erased Haskell-like
search term, then reconstructed in \Leant{}'s renderer by fitting generated
syntax back against the original Lean fragment.  This works, and Lean
verification makes it safe, but it creates a growing \emph{semantic
recovery tax}: provider metadata must be coupled to individual occurrences;
forall visibility and sort domains need parallel side channels; result types
are rediscovered from term arguments; and canonical family identity needs
special handling at every stage.

The central recommendation of this report is therefore not simply to raise the
six-binder bound, add a sextic occurrence frontier, or teach another structural
constructor to the current adapters.  The next architectural layer should be a
frontend-neutral, typed candidate graph carrying stable occurrence identities,
expected types, source-provided certificates, explicit residual obligations,
and a structured account of projection loss.  \Djex{} can continue to perform
bounded structural search over a compact logical projection; \Leant{} can keep
Lean expressions, universes, binder visibility, type-class evidence, and
indices opaque where appropriate; and the two can communicate through typed
skeletons instead of post-hoc textual repair.

On that foundation, the most promising functional extensions are: an adaptive
rather than fixed-degree rank-N plan search; a shared hash-consed subgoal and
candidate DAG; generic correlated-substitution certificates; a frontend-declared
algebra of structural families; first-class evidence and equality obligations;
recursor-skeleton synthesis for recursive data; a staged dependent-index
fragment; semantic provider retrieval at Mathlib scale; and genuinely
cooperative scheduling between \Djinn{} and \Exference{}.  The report gives
concrete data types, algorithms, invariants, migration steps, tests, and a
priority-ordered roadmap designed to preserve the projects' current fail-closed
and independently verified posture.
\end{abstract}

\tableofcontents
\clearpage
\pagenumbering{arabic}

\section{Scope, method, and terminology}

This report reviews the current source trees, the public architecture and
proposal documents, and the recent commit histories of \RepoDjex{} and
\RepoLeant{}.  The emphasis is not on listing every commit.  It is on finding
repeated causal patterns across the two repositories: what class of Lean goal
or provider exposed a gap; which repository supplied the general fix; what
information had to cross the boundary; which conservative guard was added; and
what that sequence implies for the next abstraction.

The review snapshot is the pair of main-branch heads shown on the title page.
At that point, \Leant{} pins \Djex{} exactly at
\headsha{f50b5eac86726c7530d95e48c27101b7c0fa25c9}.  This is important: the two
histories are not merely conceptually related.  The \Leant{} history records a
series of exact gitlink advances to freshly reviewed \Djex{} capabilities,
followed by Lean 4.31 transcripts that exercise those capabilities through the
complete serializer--engine--renderer--elaborator path.

Throughout the report:

\begin{itemize}
  \item \term{search type} means the type representation consumed by one of
    the synthesis engines after frontend translation;
  \item \term{source type} means the richer Haskell or Lean type known before
    projection;
  \item \term{candidate} means a generated term or term skeleton together with
    whatever evidence and obligations are needed to interpret it;
  \item \term{positive soundness} means that every displayed term is accepted
    by the source language's checker (GHC in \Djex{} integration tests, the Lean
    elaborator and kernel in \Leant{});
  \item \term{negative completeness} means that an exhausted search justifies
    a non-inhabitation claim for a stated fragment, rather than merely failing
    to find a term under finite bounds;
  \item \term{projection loss} means source structure deliberately hidden from
    the engine: an opaque dependent atom, an unopened quantifier occurrence, an
    omitted eliminator, a truncated provider inventory, or similar information.
\end{itemize}

The report deliberately distinguishes \emph{source-language semantics} from
\emph{search semantics}.  The goal is not to turn \Djex{} into a Lean
elaborator, nor to make \Leant{} duplicate both synthesis engines.  The goal is
to choose a stronger interface at which each side can retain what it uniquely
knows.

\section{Current architecture at the reviewed heads}

\subsection{Djex: two engines behind a checked neutral foundation}

\Djex{} combines \Djinn{}, a complete terminating LJT-style search for its
admitted intuitionistic fragment, and \Exference{}, a ranked heuristic search
with explicit resource controls and type-class evidence resolution.  The
merger is no longer just a shared executable.  The package exposes a
parser-independent synthesis foundation containing validated names, kinds,
types, constraints, declarations, environments, generated expressions,
queries, diagnostics, and search batches.

Several current public types are especially relevant to future work.
The shared source type is structurally richer than the original engines:

\begin{lstlisting}[language=Haskell,caption={Current shared source-type shape (abridged).}]
data Type variable
  = TypeVariable variable
  | TypeConstructor Name
  | TypeApplication (Type variable) (Type variable)
  | FunctionType (Type variable) (Type variable)
  | TupleType Boxity [Type variable]
  | ForallType
      [variable]
      [Constraint (Type variable)]
      (Type variable)
\end{lstlisting}

It distinguishes flexible inference variables from rigid skolems; provides
capture-avoiding substitution and binder normalization; preserves explicit
applications; and gives saturated functions and tuples canonical structural
forms while still exposing a constructor-application view for higher-kinded
unification.

The shared generated-term syntax has also grown beyond a first-order lambda
calculus.  It contains patterns, lambdas, applications, visible type
applications, tuples, holes, lets, and cases.  A complete application spine can
interleave term and visible-type arguments in source order.  Specified visible
type arguments are closed and alpha-normalized rather than being raw strings.

Most importantly for extensibility, a public candidate already has a natural
place for information beyond the rendered term:

\begin{lstlisting}[language=Haskell,caption={The existing candidate envelope is an opening for a richer protocol.}]
data Candidate ty details output = Candidate
  { candidateOutput              :: output
  , candidateResidualConstraints :: [Constraint ty]
  , candidateDetails             :: details
  }
\end{lstlisting}

The current \code{details} parameter is backend-owned and the residual
constraints are type-class-shaped.  The proposed architecture in this report
extends, rather than discards, that idea: make residual obligations more
general and give candidate details a stable neutral core that frontends and
both engines can share.

The provider-instantiation boundary is another major asset.  It has evolved
from independent scalar candidates into exact, ordered, correlated vectors,
and then into kinded vectors:

\begin{itemize}
  \item \code{ProviderInstantiationCandidate} contributes one provider-local
    type to a pool;
  \item \code{ProviderInstantiationAssignment} preserves a complete leading
    binder assignment in source order;
  \item \code{KindedProviderInstantiationAssignment} pairs each selected type
    with its exact ground kind.
\end{itemize}

This is already a small certificate language: an external frontend asserts a
provider identity and a correlated substitution; the checked backend validates
session identity, arity, kinds, closure, and structural fidelity before using
it.  Section~\ref{sec:certificates} argues that this should become the template
for all source-provided semantic evidence.

\subsection{Leant: semantic serializer, search adapter, and independent checker}

\Leant{} is a Lean REPL whose \code{:synth} pipeline links \Djex{} in-process.
The architecture intentionally keeps the engine boundary narrow:
\code{Leant.Synth.Engine} is the only synthesis module importing \Djex{}, while
\code{Fragment}, \code{Render}, and the backend worker own Lean-specific
translation and verification.

The input fragment is substantially richer than a simple propositional
encoding.  It distinguishes:

\begin{itemize}
  \item arrows, products, sums, top, and bottom;
  \item explicit or implicit sort quantifiers and erased instance binders;
  \item exact contextual class applications with ground-kinded arguments;
  \item bound and nominal proper-type applications;
  \item parameterized finite inductive families;
  \item parameterized recursive families with bounded constructor metadata;
  \item safe and unsafe opaque atoms; and
  \item a depth sentinel that makes truncation explicit.
\end{itemize}

The serializer runs inside Lean, after elaboration, and therefore has access to
information no Haskell parser or string renderer can reliably reconstruct:
exact constants, binder visibility, whether a domain is \code{Prop},
\code{Type}, or a general \code{Sort}, active instance heads, constructor
inventories, recursive occurrences, and definitional reduction through the
current environment.

The output path is just as important.  \Leant{} does not trust an engine term.
It fits the generated expression against the original fragment, restores exact
Lean names and binder positions, enumerates bounded source spellings where
semantic information is genuinely ambiguous, and sends each surviving term
back to Lean for elaboration against the exact goal.  A candidate is displayed
only after that check succeeds.

This independent checker is what permits aggressive positive search without
compromising correctness.  It does \emph{not}, however, make arbitrary
semantic erasure free.  Every erased fact that is later necessary for an
otherwise valid term becomes either a dropped candidate or another
renderer-side reconstruction rule.  The size and recent history of
\code{Leant.Synth.Render} and \code{Leant.Synth.Fragment} are evidence that the
cost has become architectural rather than incidental.

\subsection{Current scheduling and honesty rules}

At the reviewed head, the main interactive policy includes:

\begin{itemize}
  \item a provider-free structural baseline;
  \item live-provider fallback, with bounded discovery and staged widths;
  \item \Djinn{}, \Exference{}, or a combined portfolio;
  \item at most five displayed terms, twelve fresh candidate groups tried per
    standalone engine, and a combined frontier preserving twelve opportunities
    for each engine;
  \item an internal candidate collection window of sixty;
  \item strict \Exference{} search first, followed by an allow-unused retry only
    after a strict miss; and
  \item Lean verification of every textual variant before display.
\end{itemize}

Negative verdicts follow a stricter rule than positive ones.  \Exference{}
never provides a non-inhabitation proof.  \Djinn{} exhaustion is exposed as a
sound refutation only when the logical projection is complete and the Lean
translation did not hide concrete structure in an unsafe atom.  Provider lanes
may still find a constructive term after a provider-free refutation; if they do
not, \Leant{} restores the retained refutation before considering the classical
fallback.

This separation between positive verification and negative completeness is one
of the strongest design decisions in both projects.  Future extensions should
make it more expressive, not weaken it.

\section{The co-evolution pattern in recent history}

\subsection{A condensed cross-repository timeline}

Table~\ref{tab:timeline} groups the most informative recent changes into
capability loops.  Dates are commit dates; the short hashes link to the exact
changes.

\small
\begin{longtable}{L{2.0cm} L{3.2cm} L{4.15cm} L{4.15cm}}
\caption{Representative capability loops in the coordinated history.}
\label{tab:timeline}\\
\toprule
\textbf{Date} & \textbf{Theme} & \textbf{Djex movement} & \textbf{Leant movement} \\
\midrule
\endfirsthead
\toprule
\textbf{Date} & \textbf{Theme} & \textbf{Djex movement} & \textbf{Leant movement} \\
\midrule
\endhead
\bottomrule
\endfoot

2026-08-06 & Bounded rank-N occurrence plans &
\djcommit{d728719ff2515a8a464398aaab3e789420cc86a6}{d728719f}
adds quintuple-open and quintuple-opaque tails, capped at 512 views per
orientation and exhaustive through eleven independent sites. &
\lecommit{80f123a4}{80f123a4} and \lecommit{3fcf7cfa}{3fcf7cfa}
advance the submodule and verify genuinely non-prefix ten- and eleven-site Lean
examples. \\

2026-08-09 & Query-closed and correlated instantiation &
\djcommit{227bdc52e66e0b07606b84e89a7c3516493a718b}{227bdc52}
adds closed query-subtree choices for local schemes;
\djcommit{077f767fd75814bbcd74bed795e10455935e397f}{077f767f} through
\djcommit{844e548de3b896d7d91ecc2c48750ba229560a3e}{844e548d}
add result-correlated guarded impredicative choices while preserving historical
prefixes. &
\lecommit{21c0bab6746f5b6f309afd183ebe6ad26909e44d}{21c0bab6} pins the reviewed
engine; \lecommit{e09bded3de507c2b28bb243421cdfe07371aaa1d}{e09bded3} and
\lecommit{15fb0962ca7cf99c6df6a8c0c0776d79febe7ac9}{15fb0962}
repair explicit provider heads and nested fitting; \lecommit{c333f2c070bdb61eee6b6558deb0a827f2bd1389}{c333f2c0}
merges live Lean coverage. \\

2026-08-10 & Scalar global aggregate elimination &
\djcommit{db25170b2aed13d2bb1249c11229cc28f809aa95}{db25170b}
lets \Exference{} derive boxed-tuple eliminators from all checked type sources
and deconstruct a scalar global once. &
\lecommit{690277ca282b8c196b6a8c34777e8689539cd582}{690277ca} carries exact
provider result fragments through case and let;
\lecommit{9893fef55dc47a529c4b4d547d6e18268e1d71f5}{9893fef5} normalizes the
engine's intrinsic pair patterns; \lecommit{65b93bcc356fcdea30478dd127eb32ef8db61e1e}{65b93bcc}
merges the checked path. \\

2026-08-10 & Provider-result specialization &
\djcommit{e10f7f018facef4445fc774118b39d716b36c02e}{e10f7f01} and
\djcommit{d918334a72532fd1bb3967a6fc64cfa0eabcc4af}{d918334a}
provide one strict, lazy, mixed term/type application-spine constructor. &
\lecommit{892117dcc9756f3ede4b8500665892781e679fb7}{892117dc} through
\lecommit{9189d509ca5dc6a377c02060beb79f8965cd9879}{9189d509}
infer provider-result substitutions from exact term evidence, thread them
through full higher-order applications, and preserve trailing quantified
structure. \\

2026-08-10 & Exact contextual assignments and fidelity &
\djcommit{03aa3820284f9117011910b7eb931568ef4255cd}{03aa3820}
accepts closed contextual polytypes in exact vectors but marks selected types
through nested structures; \djcommit{103fdd391aaa350c09f07b682d6fd70a67a44989}{103fdd39}
requires every marker-bearing sibling field to retain the selected identity. &
\lecommit{2979dbe2fe25915cdaf497dca96b934a85c6b4c1}{2979dbe2}
introduces an exact-assignment-only nominal Lean class context and private
methodless engine classes; \lecommit{647f3940e3fd7859f02fde955f30bc75d4fb08a6}{647f3940}
pins the strengthened fidelity boundary. \\

2026-08-10 & Occurrence-local rendering metadata &
The canonical visible type is intentionally language-neutral; \Djex{} keeps
specified visible arguments alpha-normalized and bounded. &
\lecommit{888e2b402529089e606a55c8374d9bbf1a7568d4}{888e2b40}
retains forall visibility and sort-domain tags;
\lecommit{fe4b2932d5ff356ac5792e26f698444d1b99fe81}{fe4b2932}
couples one metadata choice to each provider occurrence so repeated uses can
select independently. \\

2026-08-10 & Higher-kinded contextual evidence &
\djcommit{c5afc68a1b49adfb66f25d67c49402487cf9697d}{c5afc68a}
preserves explicit class-parameter kinds in checked shared environments;
\djcommit{0ea49f9b7efd995c40af44fe9a2ff27de2f177bb}{0ea49f9b}
checks a rank-N assignment constrained by a type constructor. &
\lecommit{1668c956b2488bfeb4af2659314a0317a8e97fb5}{1668c956}
serializes residual kind arity, reconstructs private kinded classes, plans a
nominal head at total arity, and rejects conflicts source-locally;
\lecommit{d3662c3162b7f5014520f02b539a5ab1a657ef02}{d3662c31}
merges Lean 4.31 verification. \\

2026-08-10 & Six-binder common boundary &
\djcommit{66955da710951fd525f7ad874d133132ffcd2cee}{66955da7}
raises the single public provider/instantiation width from five to six so
\Djinn{} and \Exference{} advance together; seven remains a productive checked
rejection. &
\lecommit{d04c1a8fb9971cd7eec07fc873016627ebb1e928}{d04c1a8f}
keeps adjacent Lean forall nodes as one six-binder scheme and derives producer,
parser, engine, and renderer bounds from \Djex{}; \lecommit{de4f562c77ffa339f97670058dde693119d76345}{de4f562c}
merges all-engine and live tests. \\

2026-08-10 & Canonical structural higher-kinded identity &
\djcommit{5a550eb409c37c00256a41b14bad5a0d3792fa68}{5a550eb4}
retains the bare boxed pair constructor through \Djinn{} exact substitutions;
\djcommit{f50b5eac86726c7530d95e48c27101b7c0fa25c9}{f50b5eac}
merges both-engine and GHC-checked support. &
\lecommit{0dbe08a7819f44c872bf091685f6b30f170f0133}{0dbe08a7}
separates canonical nominal context evidence from legacy fragments and maps
exact \code{Prod}/\code{Sum} heads to their structural engine identities;
\lecommit{0160237e4bdf6e342a55f5a86ce5d5292728024f}{0160237e}
merges the Lean-verified result. \\
\end{longtable}
\normalsize

\subsection{The recurring five-step capability loop}

The timeline shows a stable development method:

\begin{enumerate}
  \item \textbf{A precise source-language counterexample appears.}
    The useful failures are rarely abstract.  They are types such as a provider
    whose visible type argument is itself quantified, a provider returning a
    tuple containing a rank-N field, a class argument whose kind is
    \code{Type -> Type}, or a canonical \code{Prod} assignment consumed later
    in saturated form.

  \item \textbf{Leant preserves just enough semantic evidence to state the
    failure faithfully.}
    Exact names, binder visibility, sort domains, instance-head correlations,
    source fragments, and constructor inventories are kept rather than guessed
    from pretty text.

  \item \textbf{Djex generalizes the missing search or boundary invariant.}
    Examples include exact ordered assignment vectors, shared kinded classes,
    a strict mixed application spine, bounded closed/correlated instantiation,
    scalar aggregate elimination, and structural assignment-fidelity checks.

  \item \textbf{Leant advances the submodule and adapts only the remaining
    Lean-specific semantics.}
    The renderer restores source names and syntax; the serializer supplies
    metadata that cannot be language-neutral; the engine module remains the
    narrow adapter.

  \item \textbf{The exact source-language case is independently checked.}
    \Djex{} uses public-facade and GHC compilation regressions.  \Leant{} uses
    pure boundary tests and Lean 4.31 transcripts.  Malformed, over-wide,
    phantom, dictionary-dependent, and unsupported controls are generally
    added beside the positive case.
\end{enumerate}

\begin{finding}{Finding 1: the synergy is not accidental coupling; it is a productive division of semantic authority.}
\Djex{} is best placed to own finite search, alpha-safe neutral syntax,
correlated substitutions, common bounds, and engine-independent invariants.
\Leant{} is best placed to own elaborated Lean identities, universes, binder
roles, instance provenance, constructor/recursor metadata, and final kernel
checking.  The future architecture should formalize this division rather than
trying to collapse it.
\end{finding}

\subsection{What the same history warns about}

The loop is productive, but the shape of the fixes has changed.  Early work
added logical rules or widened a clearly bounded family.  The most recent work
is dominated by keeping semantic side channels synchronized across:

\begin{itemize}
  \item provider discovery and exact assignment parsing;
  \item query-wide family planning;
  \item the \Djinn{} and \Exference{} projections;
  \item generated visible type applications;
  \item occurrence-local fitting;
  \item Lean binder-domain rendering; and
  \item final elaboration.
\end{itemize}

The \code{Prod}/\code{Sum} higher-kinded change is illustrative.  A saturated
product is structural in both engines; a bare \code{Prod} is a higher-kinded
source constant; an exact assignment may carry the bare or partially applied
head; a contextual class may mention it; provider planning must not close it as
an ordinary proper-type atom; the engine must map it back to the same structural
identity used by saturated products; and the renderer must retain the exact
Lean source choice.  The implementation is correct and conservative, but each
new family with analogous dual structural/nominal behavior would repeat much
of that path.

Likewise, provider-result specialization now uses exact term arguments to infer
fragment substitutions, re-applies them before inspecting each residual
constructor, fits later arguments against newly specialized domains, and
recurses through nested provider applications.  This is valuable functionality.
Architecturally, however, it is type propagation over an untyped generated
expression---that is, a small bidirectional elaborator implemented after
search because the candidate edge did not carry the information earlier.

\begin{finding}{Finding 2: the dominant scaling problem is now semantic erasure followed by source-specific recovery.}
The present boundary is no longer too weak because it lacks another expression
constructor.  It is too weak because a generated expression does not carry a
stable typed derivation explaining how each application, elimination, and
visible specialization was justified.  Leant therefore has to reconstruct that
explanation from the original fragment and occurrence metadata.
\end{finding}

\section{Capability envelope and the limits that now matter}

\subsection{What is already unusually strong}

It is easy to understate how far the combined system has moved beyond the
classic \Djinn{} and \Exference{} projects.  At the reviewed heads, the joint
capability includes all of the following.

\paragraph{A checked, reusable neutral foundation.}
Names, types, kinds, constraints, declarations, environments, generated terms,
queries, diagnostics, and search outcomes cross explicit validation edges.
Both engines are available through one package and one public facade, while
retaining independent search semantics.

\paragraph{Positive and negative search with honest evidence classes.}
\Djinn{} supplies complete search and refutation within admitted projections;
\Exference{} supplies ranked positive search under budgets.  Leant preserves
that distinction and adds a source-language verification layer.

\paragraph{Bounded higher-rank and guarded impredicative synthesis.}
Positive forall occurrences can be opened according to a deterministic plan
family.  Hypothesis-side schemes, loaded schemes, query-local schemes, and
provider schemes can be instantiated from carefully delimited vocabularies.
Quantified choices are admitted only when supplied by the checked query or by
explicit external evidence; the engines do not invent arbitrary polytypes.

\paragraph{Exact provider evidence.}
A richer frontend can preserve the correlation among all leading type binders
selected by one instance head, including heterogeneous kinds.  Evidence stays
provider-local and is checked at the session boundary rather than being
flattened into a global Cartesian pool.

\paragraph{Nominal and structural inductive identity.}
Compatible finite families can expose constructor introduction and elimination.
Compatible recursive families have bounded positive construction in \Djinn{}
and one-layer elimination in \Exference{}.  Query-wide planning prevents
separate occurrences of the same exact family from drifting into unrelated
engine types.

\paragraph{Live source environments.}
\Leant{} discovers a bounded, goal-relevant provider inventory from the actual
Lean environment, preserves exact global names, accepts project ratings, caches
semantic inventories, and verifies every use in the current session.

\paragraph{Source-faithful rendering of difficult applications.}
Visible type applications, mixed explicit and implicit forall binders, exact
sort-domain alternatives, higher-kinded assignments, contextual rank-N
arguments, and provider-result specialization all survive the round trip often
enough to solve examples that a purely erased Haskell view could not express.

These capabilities should be treated as assets and compatibility constraints.
A future redesign should not discard the deterministic prefixes, independent
verification, fail-closed evidence checks, or public bounds that made the
current system trustworthy.

\subsection{Current finite bounds}

Table~\ref{tab:bounds} collects the most relevant finite limits at the reviewed
heads.  Some are public \Djex{} constants; some are \Leant{} policy values; and
some are documented search-family bounds.  Their exact values matter less than
the fact that they are explicit, productive, and tested on malformed or cyclic
caller-built inputs.

\begin{table}[H]
\centering
\small
\caption{Selected finite bounds at the review snapshot.}
\label{tab:bounds}
\begin{tabularx}{\textwidth}{L{4.35cm} C{1.8cm} Y}
\toprule
\textbf{Bound} & \textbf{Value} & \textbf{Role} \\
\midrule
Provider instantiation candidates per checked query & 32 &
Bounds provider-local scalar associations before search. \\
Exact provider assignments per checked query & 32 &
Bounds complete correlated vectors accepted by public runners. \\
Leading arguments in one exact assignment & 6 &
Single common eligibility boundary for current \Djinn{} and \Exference{}
instantiation routes; seven is a fail-closed control. \\
Ground-kind constructor nodes per argument & 129 &
Admits the complete right-associated \code{Type -> ... -> Type} shape through
arity 64 while productively rejecting cyclic or excessive values. \\
Active Lean instance heads inspected per provider & 32 &
Bounds frontend discovery in Lean resolver order. \\
Retained alpha-distinct vectors per live provider & 16 &
Preserves correlated source choices without unbounded metadata growth. \\
Live providers serialized by Leant & 80 &
Bounds one discovered inventory before staged search. \\
Semantic provider-cache entries & 12 &
Generation-aware LRU keyed by canonical query/provider-world information. \\
Displayed candidates & 5 &
Interactive result frontier. \\
Fresh groups tried per standalone engine & 12 &
Backend verification budget after source-local deduplication. \\
Combined fresh-group frontier & 24 &
Preserves twelve opportunities for each engine under the current merge order. \\
Engine candidate collection window & 60 &
Bounds \Djinn{} sorted collection and the shared ranking input. \\
Rank-N quintic plan views & 512 per orientation &
Covers all relevant ten- and eleven-site fifth-order selections; a twelve-site
six/six split remains the documented gap. \\
Instantiation tuple attempts & 512 per scheme family &
Bounds fair candidate-tuple production. \\
Retained axioms & 16 per scheme; 64 per family &
Bounds positive-only generated specialization premises. \\
\bottomrule
\end{tabularx}
\end{table}

The right future direction is not to remove these bounds.  It is to make each
bound measure a semantic resource rather than a historical implementation
dimension.  For example, ``at most 1,024 distinct compiled formula views'' is a
more general budget than ``quintuple-open plus quintuple-opaque''; ``at most 32
validated substitution certificates'' is more general than a provider-only
assignment vector; and ``at most 60 semantic candidate fingerprints'' is more
useful than sixty derivations that later render to duplicates.

\subsection{Five boundaries that block substantially more automation}

\subsubsection{The generated term is mostly untyped at the neutral edge}

The search engines know types internally, and \Leant{} knows the exact expected
Lean fragment, but the shared generated expression is an unannotated syntax
tree.  Once a candidate crosses that edge, local type information, rule
provenance, instantiation substitutions, and branch refinements are represented
only indirectly through global maps and a re-fitting traversal.

This is the root cause behind several recent fixes:

\begin{itemize}
  \item recovering a provider's specialized result from exact term arguments;
  \item carrying rank-N field types through case and let patterns;
  \item deciding where explicit forall placeholders must be woven;
  \item coupling a render-metadata choice to a specific provider occurrence;
  \item preserving trailing forall structure during higher-order application;
  \item distinguishing a canonical structural head from a merely similar
    nominal source fragment.
\end{itemize}

\subsubsection{Quantifier occurrence planning is encoded by degree}

The positive rank-N plan family has progressed from singleton choices through
pairs, triples, quadruples, and capped quintuples.  The implementation is
carefully prefix-preserving and the coverage results are strong: the quintic
family is exhaustive through eleven independent sites.  But the next gap is
already known: a twelve-site six-open/six-opaque split.  Adding a sextic family
would close that case and expose the next central antichain.

This is a sign that the representation of the search problem should change.
The semantic object is a finite partially ordered set of occurrence decisions
with ancestor prerequisites and formula equivalences.  ``Degree five'' is one
enumeration strategy for that object, not an enduring abstraction.

\subsubsection{Source evidence is divided among specialized side channels}

Current evidence channels include scalar provider candidates, exact assignment
vectors, kinded vectors, provider binder names, exact forall-domain tags,
contextual class fragments, nominal higher-kinded heads, constructor maps,
type-name maps, and source fragments retained for fitting.  Each addition has a
reasonable local purpose.  Together, they describe a more general object: a
source-certified substitution or family fact with a stable provenance and a
specified lifetime.

Until those channels share one certificate model, every new semantic fact must
answer the same questions separately: how is it bounded, alpha-normalized,
associated with a provider occurrence, invalidated with the session, projected
into both engines, rendered in the source language, and rejected without
poisoning unrelated evidence?

\subsubsection{Dependent structure is transportable but not refinable}

The current opaque-atom discipline is exactly right for many dependent goals.
If two alpha-equivalent dependent propositions merely need to be transported,
the engine can treat them as one atom and synthesize the higher-order plumbing.
But constructor analysis of indexed families, equality transport, dependent
pairs, subtypes, and motives remain beyond the structural projection.

A binary choice between ``fully opaque'' and ``fully implement dependent type
theory in the engine'' would be a mistake.  There is a valuable middle layer:
search a structural term skeleton while retaining source-owned equality,
index, universe, coercion, and type-class obligations as explicit holes.

\subsubsection{Environment search is bounded by inventory width rather than
semantic reachability}

The provider-free baseline and geometric widths of 1, 4, 16, and the complete
bounded inventory protect structural search from a large environment.  They do
not answer the harder retrieval question: which declarations are jointly
useful for a type that requires two or three compositions, perhaps through
higher-ranked or constrained intermediates?

A persistent Mathlib-scale inventory is already identified as future work in
the project proposal.  The key extension is not merely persistence.  It is a
semantic index and a backward relevance expansion driven by result types,
argument requirements, family identities, constraints, and learned successful
compositions.

\section{Central diagnosis: the semantic recovery tax}
\label{sec:recovery-tax}

\subsection{A concrete model of the current round trip}

A simplified current path for a difficult provider candidate is:

\begin{enumerate}
  \item Lean elaborates the provider and goal.
  \item The serializer projects them to \code{Frag}, preserving selected exact
    metadata in parallel records.
  \item \code{Engine} lowers \code{Frag} into shared \Djex{} types and private
    declarations, often erasing Lean binder roles or replacing exact families
    with collision-free engine names.
  \item \Djinn{} or \Exference{} searches and emits a neutral
    \code{Expression} with visible type applications where the engine retained
    them.
  \item \code{Render} reconstructs a local domain environment, strips synthetic
    premises, aliases provider occurrences, chooses source metadata,
    specializes provider results from exact arguments, fits nested term
    arguments, introduces missing binders, normalizes patterns, and enumerates
    textual variants.
  \item Lean elaborates every variant against the original goal; failures are
    silently discarded.
\end{enumerate}

The path is safe because step 6 is authoritative.  Its incompleteness comes
from steps 3--5: once information is erased, the renderer can recover it only
when the original fragment and the generated syntax contain enough clues.

\subsection{Why verification alone is not a substitute for typed candidates}

One might argue that Lean is the checker, so the engine can emit arbitrary
terms and let elaboration decide.  That is correct for soundness but poor for
search completeness and cost.

\begin{itemize}
  \item \textbf{Rejected candidates consume the expensive resource.}
    The engine's pure search is often cheap relative to a backend round trip.
    Missing a type argument or choosing the wrong forall-domain spelling turns
    one semantic candidate into many elaboration attempts.

  \item \textbf{Elaboration failure is currently unstructured feedback.}
    A rejected spelling does not tell the engine whether the problem was a
    universe, binder visibility, an unsolved instance, a stale provider result,
    an index mismatch, or a genuinely wrong proof skeleton.

  \item \textbf{Untyped deduplication is weak.}
    Distinct derivations may normalize to the same source term; source variants
    of one semantic candidate may occupy multiple slots; and two terms with the
    same printed skeleton can differ in exact instantiation provenance.

  \item \textbf{Negative evidence cannot use verification failures.}
    A finite list of rejected candidates says nothing about inhabitation.  The
    projection and search completeness must still be tracked independently.
\end{itemize}

\subsection{The desired replacement}

The replacement is not a fully source-specific core language.  It is a
\term{typed synthesis skeleton}:

\begin{itemize}
  \item every term edge has an engine-visible expected or inferred neutral
    type;
  \item every local and global use has a stable occurrence identity;
  \item every explicit specialization records the substitution certificate
    that justified it;
  \item every elimination records the family descriptor and branch-local field
    types it exposes;
  \item facts that remain source-specific become named residual obligations;
  \item source rendering metadata is attached to the exact node whose choice it
    affects; and
  \item the candidate carries a derivation/provenance DAG sufficient for a
    small independent neutral checker.
\end{itemize}

\begin{recommendation}{Recommendation 1: introduce a typed candidate graph before adding another major type-language feature.}
This is the highest-leverage architectural change.  It directly simplifies
provider-result fitting, occurrence-local metadata, structural elimination,
semantic deduplication, elaboration diagnostics, and future dependent or
recursive skeletons.  It can be added incrementally beside the existing
\code{Expression} output and projected back to that syntax for compatibility.
\end{recommendation}

\section{Proposed target architecture}
\label{sec:architecture}

\subsection{Four layers with explicit contracts}

The target architecture should make four layers visible.

\begin{description}
  \item[Layer A: source semantic package.]
    Produced by the frontend after source elaboration.  It contains a neutral
    type skeleton, opaque source atoms, exact family descriptors, provider
    certificates, binder roles, and bounded rendering metadata.  In \Leant{},
    this package is produced by Lean metaprograms.  In a Haskell frontend, it
    is produced from the checked source environment.

  \item[Layer B: logical/search projection.]
    A compact representation tailored to \Djinn{} and \Exference{}.  It may
    erase universes, indices, dictionaries, and some quantifier occurrences,
    but it records every erasure as a structured projection-loss fact and
    retains stable links back to Layer A.

  \item[Layer C: typed candidate and derivation graph.]
    Search produces typed skeletons, not only surface expressions.  Nodes
    refer to source/global/family identities through stable handles.  Edges
    carry expected types, substitutions, and residual obligations.  Both
    engines can project this graph to the current generated syntax.

  \item[Layer D: source completion and verification.]
    The frontend renders or directly builds a source AST, solves residual
    obligations using its elaborator, verifies the candidate, and optionally
    returns structured failure information to a bounded refinement loop.
\end{description}

The crucial rule is that no layer is asked to rediscover information owned by a
previous layer.  If a source fact is intentionally hidden from search, the
candidate carries a reference or obligation, not an invitation for the
renderer to infer it from text.

\subsection{A frontend-neutral typed type skeleton}

The existing shared \code{Type} should remain the Haskell-facing source type
and the common type for the current engines.  A new adjacent representation can
express the additional distinctions required by richer frontends without
forcing every engine to understand them.

One possible shape is:

\begin{lstlisting}[language=Haskell,caption={Sketch of a richer neutral source-semantic type.}]
data SortClass
  = PropositionSort
  | DataSort
  | GeneralSort SortToken

data BinderRole
  = TypeBinder
  | TermBinder
  | InstanceBinder

data Visibility
  = Explicit
  | Implicit
  | StrictImplicit
  | InstanceImplicit

data SourceType variable atom
  = STVariable variable
  | STNominal FamilyId [SourceType variable atom]
  | STFunction (SourceType variable atom)
               (SourceType variable atom)
  | STPi (Binder variable atom)
         (SourceType variable atom)
  | STProduct ProductFlavor [SourceType variable atom]
  | STSum SumFlavor [SourceType variable atom]
  | STIndexed FamilyId
      [SourceType variable atom] -- parameters
      [atom]                     -- source-owned indices
  | STOpaque atom

data Binder variable atom = Binder
  { binderIdentity   :: BinderId
  , binderSourceName :: Maybe SourceName
  , binderRole       :: BinderRole
  , binderVisibility :: Visibility
  , binderSort       :: SortClass
  , binderDomain     :: SourceType variable atom
  }
\end{lstlisting}

This is deliberately not a complete Lean expression language.  An index term
or a universe expression may remain an opaque, hash-consed source token with
frontend-provided alpha/definitional-equality relations.  The search engines
need only the structural distinctions that affect legal rules.  A frontend can
supply more exact data without making \Djex{} depend on Lean.

The current \code{Frag} can initially be the \Leant{} implementation of
\code{SourceType}.  The migration benefit comes from moving its reusable
semantics---binder roles, exact family identities, occurrence IDs, and source
obligations---to a public neutral package consumed directly by the typed
candidate edge.

\subsection{Typed candidate nodes and residual obligations}

A candidate graph can be represented with typed node references rather than a
recursive tree.  Hash-consing then becomes natural, repeated subterms share
storage, and provenance can refer to one stable node.

\begin{lstlisting}[language=Haskell,caption={Sketch of a typed candidate envelope.}]
data TypedCandidate ty local source = TypedCandidate
  { candidateRoot        :: TermNodeId
  , candidateGraph       :: TermGraph ty local source
  , candidateObligations :: [Obligation ty source]
  , candidateDerivation  :: DerivationGraph
  , candidateLosses      :: ProjectionLossSet
  , candidateCost        :: CostVector
  , candidateFingerprint :: CandidateFingerprint
  }

data TermNode ty local source
  = LocalNode local ty
  | GlobalNode GlobalId ty source
  | LambdaNode [TypedPattern ty local] TermNodeId ty
  | ApplyNode TermNodeId TermNodeId ApplicationWitness ty
  | TypeApplyNode TermNodeId SubstitutionCertificateId ty
  | ConstructorNode ConstructorId [TermNodeId] ty
  | EliminateNode EliminatorId TermNodeId [BranchNode] ty
  | LetNode (TypedPattern ty local) TermNodeId TermNodeId ty
  | ObligationNode ObligationId ty

data Obligation ty source
  = TypeClassGoal source ty
  | EqualityGoal source ty ty
  | CoercionGoal source ty ty
  | UniverseGoal source
  | ElaborationChoice source [SourceAlternative]
  | TerminationGoal source
  | FrontendGoal source ty
\end{lstlisting}

\code{ApplicationWitness} need not contain a full proof object.  It can record
which neutral function type and substitution the engine used, enough for a
small checker to re-run the typing step.  \code{source} is an opaque frontend
payload or handle, bounded and validated when the query is sealed.

This design makes the current \code{candidateResidualConstraints} a special
case of \code{candidateObligations}.  Ordinary Haskell type-class constraints
remain directly representable.  Lean instance synthesis, equality refinement,
universe ambiguity, and coercion insertion become additional obligation kinds
which \Djex{} may carry without solving.

\subsection{Stable occurrence identities}

Recent \Leant{} work had to create collision-safe aliases for separate uses of
the same provider because two occurrences with the same canonical visible type
may need different source visibility/domain alternatives.  That is a strong
argument for occurrence identity in the neutral graph.

Every source provider use, quantified occurrence, family occurrence, and
opaque atom should have an \code{OccurrenceId} allocated by the checked
frontend package.  Search-created nodes may have separate internal IDs.  The
mapping must obey three invariants:

\begin{enumerate}
  \item IDs are stable within one sealed query and never inferred from printed
    names.
  \item Source metadata is keyed by the exact occurrence whose interpretation
    it changes.
  \item Alpha-equivalent types may share a semantic fingerprint without losing
    distinct occurrence identity.
\end{enumerate}

This removes the need to synthesize fake global names solely to index metadata,
and it allows multiple uses of one provider to make independent bounded
choices while still normalizing to the same global source identity.

\subsection{A typed, bounded source-certificate table}
\label{sec:certificates}

The current exact provider assignment is the prototype for a general source
certificate.  A useful abstraction is:

\begin{lstlisting}[language=Haskell,caption={Generalized source-certified substitution evidence.}]
data SubstitutionCertificate ty source = SubstitutionCertificate
  { certificateOwner       :: EvidenceOwner
  , certificateTelescope   :: TelescopeFingerprint
  , certificateSubstitution:: [KindedArgument ty]
  , certificatePremises    :: [Obligation ty source]
  , certificateProvenance  :: source
  , certificateFidelity    :: FidelityWitness
  , certificateRendering   :: RenderingMetadata
  }

data EvidenceOwner
  = ProviderOwner GlobalId
  | LocalSchemeOwner OccurrenceId
  | FamilyOwner FamilyId
  | FrontendRuleOwner FrontendRuleId
\end{lstlisting}

The checked query boundary validates:

\begin{itemize}
  \item owner/session identity;
  \item exact telescope length and binder order;
  \item kinds and closure;
  \item capture safety and alpha normalization;
  \item provenance-size bounds;
  \item structural fidelity where the engine relies on nominal distinctions;
  \item and source-local conflict handling.
\end{itemize}

Both engines then consume the same certificate, perhaps through different
projections.  \Leant{} retains source rendering metadata under the certificate
ID, not in a separate occurrence-guessing map.  This generalization can cover
active instance-head assignments, query-correlated choices, exact class
contexts, family equalities, and future frontend-derived coercion or recursor
facts.

\subsection{Frontend-declared family algebra}
\label{sec:family-algebra}

The system needs a single language-neutral description of a structural family.
Hard-coded recognition of functions, tuples, lists, \code{Either},
\code{Prod}, and \code{Sum} was a reasonable bootstrapping strategy.  The
recent higher-kinded structural-identity work shows why it should now be
factored.

\begin{lstlisting}[language=Haskell,caption={Sketch of a family descriptor supplied at the checked boundary.}]
data FamilyDescriptor ty source = FamilyDescriptor
  { familyId          :: FamilyId
  , familyKind        :: Kind
  , familySortClass   :: SortClass
  , familyParameters  :: [ParameterDescriptor]
  , familyIndices     :: [IndexDescriptor source]
  , familyConstructors:: [ConstructorDescriptor ty source]
  , familyEliminators :: [EliminatorDescriptor ty source]
  , familyRecursion   :: RecursionDescriptor
  , familyLaws        :: FamilyLawSet
  , familyRendering   :: source
  }

data ParameterRole
  = RepresentationalParameter
  | PhantomParameter
  | IndexDeterminingParameter
  | RecursiveParameter
\end{lstlisting}

The descriptor should state, rather than force the engines to rediscover:

\begin{itemize}
  \item exact nominal identity and total arity;
  \item parameter versus index positions;
  \item the sort of the result;
  \item constructor field schemas after parameter substitution;
  \item whether elimination is finite, one-layer recursive, recursor-based, or
    unavailable;
  \item recursive SCC membership and strictly positive positions;
  \item whether a higher-kinded unsaturated head shares identity with an
    existing structural form;
  \item phantom/faithful parameter flow;
  \item and source rendering or AST handles.
\end{itemize}

For Haskell, the descriptor can be derived from checked data declarations and
intrinsic tuple/function syntax.  For Lean, it can be derived from
\code{InductiveVal}, constructor telescopes, and recursor metadata.  The
engines still decide which subset they support.  Unsupported laws are recorded
as projection loss, not silently approximated.

This one abstraction would subsume a surprising amount of recent special-case
work: query-wide family identity, total-arity planning, constructor maps,
recursive occurrence descriptions, \code{Prod}/\code{Sum} structural identity,
phantom-fidelity analysis, and source-specific constructor rendering.

\subsection{Capability-indexed projection and evidence}

The current \code{SynthRefuted Bool} is simple and safe, but the Boolean is
becoming too coarse.  A future outcome should explain what was complete and
what was not.

\begin{lstlisting}[language=Haskell,caption={Structured projection loss and negative evidence.}]
data ProjectionLoss source
  = OpaqueDependentAtom OccurrenceId source
  | BoundedQuantifierPlan [OccurrenceId]
  | UnsupportedFamilyElimination FamilyId
  | RecursiveDepthBound FamilyId Int
  | ProviderInventoryTruncated Int
  | ProviderEvidenceTruncated GlobalId
  | SearchBudgetReached SearchBudgetKind
  | FrontendDepthReached OccurrenceId

data NegativeEvidence source
  = RefutedInFragment CapabilityFingerprint RefutationCertificate
  | Inconclusive [ProjectionLoss source]
\end{lstlisting}

A refutation remains user-facing as a concise sentence, but internally it is
indexed by a capability fingerprint: which connectives were exact, which
families exposed complete constructor rules, whether every relevant forall
occurrence was covered, whether provider premises were intentionally excluded,
and whether any source atom hid analyzable structure.

This yields two benefits.  First, new positive-only extensions can be added
without threading another ad-hoc completeness flag.  Second, a later frontend
or engine can compare capability fingerprints and decide whether one
refutation strictly dominates another.

\begin{recommendation}{Recommendation 2: generalize exact provider evidence and family metadata into checked source certificates.}
Do this after the first typed-candidate slice, not before.  Certificates give a
uniform home to the information already flowing through provider assignments,
rendering maps, family plans, and Lean-side evidence.  They also preserve the
projects' strongest invariant: untrusted caller-built metadata is validated at
one finite boundary and cannot poison unrelated evidence.
\end{recommendation}

\section{Adaptive rank-N occurrence planning}
\label{sec:adaptive-plans}

\subsection{Why another fixed-degree frontier is the wrong primary move}

The current occurrence-plan history is disciplined:

\begin{itemize}
  \item singleton frontiers established the basic positive-forall choice;
  \item pair, triple, and quadruple tails expanded complete coverage while
    preserving historical results;
  \item the quintic tail alternates selections from both source edges, caps
    each orientation at 512 compiled views, and covers every ten- and
    eleven-site fifth-order choice;
  \item explicit sentinels identify the next uncovered central split.
\end{itemize}

This is good engineering.  It should remain as the deterministic compatibility
prefix.  But a sextic implementation would encode the same subset lattice one
more time.  The underlying problem can be searched directly.

\subsection{Plan lattice}

Give every positive forall occurrence a stable ID.  A plan assigns each
occurrence one of a small number of modes, initially:

\[
  \mathsf{Opaque},\qquad \mathsf{Open}.
\]

Nested occurrences impose prerequisite closure: opening a child requires
opening the ancestor path that exposes it.  Future modes may include
\(\mathsf{OpenWithCertificate}(c)\) or a source-specific obligation mode, but
the binary lattice is enough for the first implementation.

A plan is normalized by:

\begin{enumerate}
  \item ancestor closure;
  \item elimination of choices for occurrences unreachable under the plan;
  \item alpha-normalization of the resulting projected formula; and
  \item hash-consing against previously compiled formula views.
\end{enumerate}

Two distinct bit vectors that compile to the same formula consume one semantic
plan slot, not two.

\subsection{Compatibility prefix followed by best-first expansion}

The scheduler should emit the current historical sequence unchanged through
the quintic tails.  If no term or conclusive result appears, continue with a
best-first queue over normalized plans.

A practical priority can combine:

\[
\begin{aligned}
\mathsf{score}(p) ={}&
  w_1\,\mathsf{newOpenSites}(p)
 +w_2\,\mathsf{projectedFormulaSize}(p)\\
 &+w_3\,\mathsf{estimatedPremiseCount}(p)\\
 &+w_4\,\mathsf{distanceFromFailedPlans}(p)
 -w_5\,\mathsf{queryCorrelation}(p).
\end{aligned}
\]

The exact weights are policy, not correctness.  Determinism requires a stable
tie-break by source occurrence order and normalized plan fingerprint.

Candidate successors can be generated by:

\begin{itemize}
  \item flipping one boundary occurrence;
  \item adding the smallest unmet ancestor-closed set;
  \item moving toward the complement of a plan whose search was structurally
    unproductive;
  \item or satisfying a conflict set returned by the formula compiler or
    source verifier.
\end{itemize}

The budget should be stated in semantic units: maximum distinct compiled
formula views, maximum cumulative formula nodes, maximum prepared-premise
cache size, and maximum search work.  The current 1,024 quintic views provide a
reasonable starting ceiling for the adaptive tail.

\subsection{Counterexample-guided refinement without trusting failures}

Lean elaboration failures cannot justify negative evidence, but they can guide
positive refinement.  Suppose a candidate skeleton fails because one provider
use requires an explicit type application whose source forall remained opaque.
If the verifier returns a structured failure associated with occurrence
\(o\), the adaptive planner can prioritize plans opening \(o\) or select a
certificate that fixes its source instantiation.

The feedback loop is heuristic only:

\begin{itemize}
  \item a verifier rejection may increase a successor's priority;
  \item it never removes a plan needed for completeness;
  \item it never turns a bounded miss into a refutation; and
  \item all historical deterministic plans remain available independently of
    feedback quality.
\end{itemize}

Even without verifier feedback, the engine can derive useful conflict sets.
Examples include a quantified hypothesis whose body cannot match any current
goal subtree until a particular occurrence is opened, or a compiled plan in
which all newly exposed premises are duplicates of existing axioms.

\subsection{Data structures}

The first implementation does not require a sophisticated decision diagram.
A persistent bit set plus a hash map keyed by compiled formula fingerprint is
sufficient at current scales.  If plans grow to dozens of sites, a zero-suppressed
BDD is a natural representation for sparse occurrence subsets, while a normal
BDD can compact complement-related open/opaque families.  The key is to keep
that optimization below the semantic interface.

\begin{recommendation}{Recommendation 3: preserve the current quintic prefix, then replace degree-by-degree widening with an adaptive normalized plan lattice.}
This closes the known twelve-site central gap without committing the project to
a permanent sequence of sextic, septic, and higher frontiers.  It also creates
a place to use structured verification feedback and query-correlation scores
while retaining finite, deterministic search.
\end{recommendation}

\section{Shared subgoal DAG, semantic identity, and dominance}
\label{sec:dag}

\subsection{Why both engines should share more than input declarations}

The current combined mode runs \Djinn{} and \Exference{} independently and
then merges candidate groups with a carefully designed fairness schedule.  This
preserves each engine's semantics, but it duplicates work at three levels:

\begin{itemize}
  \item equivalent type substitutions and opened-forall views are prepared
    separately;
  \item structurally identical subgoals are explored under independently built
    environments; and
  \item semantically duplicate candidates are often recognized only after
    rendering or exact-spelling comparison.
\end{itemize}

The engines need not share one search algorithm to share immutable canonical
objects.  A neutral query preparation layer can intern types, substitutions,
family descriptors, local contexts, and subgoal keys.  Each engine then keeps
its own frontier and proof theory while reusing the same identities and caches.

\subsection{Canonical subgoal keys}

A subgoal cache key should include at least:

\[
\begin{aligned}
K = (&\text{goal type},\ \text{local-environment fingerprint},\\
     &\text{rigid scope},\ \text{available certificates},\\
     &\text{usage policy},\ \text{projection capability}).
\end{aligned}
\]

The local environment fingerprint must be order-sensitive where the engine's
enumeration order matters, but it can separately carry an order-insensitive
logical fingerprint for cache reuse.  This distinction mirrors the existing
oldest-first fairness work: proof identity and candidate ordering are related
but not identical concerns.

Prepared premise families, rank-N plan views, and exact provider substitutions
can all reference the same interned type nodes.  The current prepared-premise
caches in \Djinn{} are a natural starting point.

\subsection{Semantic candidate fingerprints}

A candidate fingerprint should be computed before source rendering.  It must
preserve semantic distinctions that affect typing while eliminating irrelevant
implementation history.

Recommended normalization steps are:

\begin{enumerate}
  \item alpha-normalize local binders and source-owned closed type binders;
  \item flatten and rebuild mixed application spines in canonical source order;
  \item normalize explicit constructor identities through family descriptors;
  \item perform safe beta reduction of administrative lambdas and let nodes;
  \item eta-contract only when the typed graph proves the contraction preserves
    explicit source binder obligations;
  \item canonicalize irrefutable tuple patterns and branch ordering where the
    family descriptor permits it;
  \item retain certificate IDs only through their semantic substitution and
    obligation fingerprint, not allocation order; and
  \item include residual obligations and expected result type.
\end{enumerate}

The result is not necessarily a source-language normal form.  It is a stable
identity for scheduling and deduplication.  Source renderers may still offer
multiple spellings of one fingerprint.

This change would improve several current limits immediately.  The sixty-item
collection window would contain sixty semantic terms rather than sixty backend
derivations.  The twelve-group verification frontier would not be spent on
terms already rejected under the same source semantics.  Combined mode could
merge by fingerprint before applying its \(D_1\ldots D_4, E_1\ldots E_{12},
D_5\ldots D_{12}\) fairness schedule.

\subsection{Dominance pruning}

For one semantic term, prefer a candidate that:

\begin{itemize}
  \item has fewer or weaker residual obligations;
  \item relies on a smaller projection-loss set;
  \item uses fewer providers or lower total provider penalty;
  \item has a smaller typed term graph;
  \item needs fewer source rendering alternatives; and
  \item preserves more exact source metadata.
\end{itemize}

This defines a partial order rather than one scalar score.  Keep the Pareto
frontier per semantic fingerprint.  A source frontend may accept a slightly
larger term because it has no unsolved type-class obligation, while a Haskell
frontend may prefer the opposite.  The current engine-specific details can
remain as tie-break data below this neutral dominance relation.

\subsection{Optional equality saturation}

An e-graph is not required for the first implementation, but the typed term DAG
makes it possible later.  A small, explicitly typed rewrite set could identify
administrative equivalences such as:

\begin{itemize}
  \item beta/eta and let-floating within capture-safe scopes;
  \item tuple projection after known construction;
  \item case-of-constructor simplification;
  \item reassociation of engine-generated applications; and
  \item elimination of synthetic premise lambdas after substitution.
\end{itemize}

The rewrite set must never use proof irrelevance or source-specific definitional
equality unless the frontend supplies an exact certificate.  Equality
saturation is therefore an optimization for candidate identity and simplification,
not a replacement for source checking.

\begin{recommendation}{Recommendation 4: deduplicate and rank typed semantic candidates before rendering.}
This is likely the best near-term latency improvement after the typed edge.  It
reduces expensive Lean verification, improves the meaning of existing windows,
and gives combined mode a shared unit of fairness without merging the engines'
search semantics.
\end{recommendation}

\section{First-class evidence and source obligations}
\label{sec:obligations}

\subsection{From erased dictionaries to explicit proof goals}

The present architecture deliberately avoids treating Lean instance binders as
ordinary proof power inside \Djinn{}.  A provider use can leave a class argument
implicit and let Lean's elaborator rebuild it.  Exact assignment discovery can
use active instance heads to determine type arguments, but the dictionary
itself is not transported as a general engine value.

That conservative boundary is sound and has avoided cross-language confusion.
It also excludes useful terms whose structure depends on an available class
method or whose result has residual evidence obligations.

A safe extension is not to import Lean dictionaries as executable Haskell
terms.  It is to let the candidate graph contain an opaque evidence node:

\begin{lstlisting}[language=Haskell]
EvidenceNode
  { evidenceGoal       :: SourceEvidenceGoalId
  , evidenceNeutralType:: NeutralConstraint
  , evidenceExpected   :: Type
  }
\end{lstlisting}

The frontend certifies that the goal is well-formed in the sealed environment.
The engine can carry the node, compose providers around it, and charge it a
cost.  The source completion layer asks Lean's instance synthesizer to solve it
in the exact local context.  Failure rejects the candidate; success fills the
hole with a source expression not interpreted by \Djex{}.

This supports a useful middle ground:

\begin{itemize}
  \item class-constrained provider composition;
  \item \code{Decidable} and coercion obligations;
  \item higher-kinded constraints retained inside rank-N arguments;
  \item local instance hypotheses in prove mode; and
  \item source-specific elaboration mechanisms that have no meaningful Haskell
    analogue.
\end{itemize}

\subsection{Evidence tables rather than arbitrary callbacks}

For determinism, purity, and cacheability, the preferred first design is a
bounded evidence table produced before search, not an arbitrary engine-to-Lean
callback during lazy evaluation.  A frontend request can include:

\begin{itemize}
  \item exact evidence goals known to be solvable in the current environment;
  \item the type variables or source metavariables they determine;
  \item complete correlated substitutions, where relevant;
  \item a stable provenance token and environment generation; and
  \item a cost or priority supplied by source resolution order.
\end{itemize}

Search may use only table entries.  A later bounded refinement round may ask the
frontend for additional evidence goals discovered in a candidate skeleton,
but each round seals a new immutable request.

This reproduces the strongest properties of current active-instance
assignments: isolation between attempts, exact vector correlation, finite
limits, and no mutation of the next candidate's metavariable state.

\subsection{Equality and coercion obligations}

The same mechanism can carry source equalities without teaching the engines a
full dependent conversion checker.  An \code{EqualityGoal} contains opaque
source terms and their neutral types.  Search rules may create it only through
a family descriptor or frontend certificate.  The source layer discharges it
using definitional equality, an available proof, a coercion, or an explicit
transport term.

For example, a constructor branch of an indexed family may refine an index
from \(n\) to \(0\) or \(m+1\).  The engine can solve the structural branch
body under a source refinement token while leaving the exact equality proof to
Lean.  The token is scoped to that branch and cannot escape unless the
candidate carries a corresponding equality obligation.

\subsection{Obligation-aware cost}

Residual obligations should be visible to ranking.  A reasonable default cost
order is:

\[
\begin{aligned}
\text{no obligation}
&< \text{known table evidence}
 < \text{ordinary instance goal}\\
&< \text{definitional equality}
 < \text{explicit equality transport}\\
&< \text{open frontend hole}.
\end{aligned}
\]

The order is policy and frontend-dependent.  The important invariant is that a
term with unresolved obligations is never presented as complete before the
source layer fills and checks them.

\section{Recursive synthesis through recursor skeletons}
\label{sec:recursion}

\subsection{Current conservative recursive support}

The current split is principled:

\begin{itemize}
  \item \Djinn{} can construct bounded positive recursive values, limited by
    SCC-aware rules and an independent-SCC cap;
  \item \Exference{} can perform one finite constructor match and treats
    recursive fields as ordinary locals without recursively eliminating them;
  \item neither engine invents an unbounded recursive definition or performs
    induction; and
  \item Lean or GHC remains the final checker of generated terms.
\end{itemize}

This covers finite constructor terms and useful one-layer projections while
keeping termination obvious.  The next step should preserve that posture.

\subsection{Recursors as ordinary certified providers}

In Lean, each inductive family already has recursors and cases eliminators with
checked termination.  Rather than adding general recursion to the engine,
serialize a bounded \term{recursor descriptor} as part of the family algebra.
The descriptor identifies:

\begin{itemize}
  \item the motive parameters and indices;
  \item the major premise;
  \item one branch telescope per constructor;
  \item which branch fields receive recursive hypotheses;
  \item universe/sort restrictions; and
  \item a source handle for rendering the exact recursor application.
\end{itemize}

Search can then synthesize a recursor skeleton:

\[
  \mathsf{recursor}\;\mathsf{motive}\;
  \mathsf{branch}_1\;\cdots\;\mathsf{branch}_k\;\mathsf{major}.
\]

Each branch is a typed subgoal.  Recursive hypotheses are ordinary local
providers whose types come from the descriptor.  \Djinn{} or \Exference{} can
solve branch plumbing without ever constructing a self-reference.

For Haskell, the analogous descriptor can point at an existing trusted fold or
eliminator such as \code{foldr}, rather than requiring the neutral term syntax
to express \code{let rec}.  A future Haskell-specific frontend may add checked
structural recursion separately.

\subsection{When to propose a recursor}

Unrestricted recursor insertion can explode the search space.  Gate it by
source and goal relevance:

\begin{enumerate}
  \item a local or provider value has an exact recursive family at its head;
  \item the target type depends on, or is reachable from, fields exposed by one
    constructor layer;
  \item one-layer case analysis failed or leaves a branch goal whose type
    matches the recursive result at a structurally smaller field;
  \item a recursor descriptor is complete and source-certified; and
  \item the per-query recursor budget has not been spent.
\end{enumerate}

Start with non-dependent catamorphism-like recursors for \code{Nat},
\code{List}, and simple user families.  The branch result type must be uniform
and first-order in the initial slice.  Later, the same typed skeleton can carry
motive and index obligations.

\subsection{Bounded recursion policies}

Useful finite controls include:

\begin{itemize}
  \item at most one recursor skeleton on a term path initially;
  \item at most one recursive family SCC per skeleton;
  \item a cap on branch subgoals and recursive-hypothesis uses;
  \item no synthesized nested recursor until all one-recursor candidates have
    been tried;
  \item a size penalty dominating ordinary application and one-layer cases;
  \item and no negative completeness claim for projections that omit an
    available recursor.
\end{itemize}

These bounds are semantic and can be widened independently of the source
language.  Lean's termination checker remains authoritative because the
neutral engine never generates an unchecked recursive binding.

\begin{recommendation}{Recommendation 5: add certified recursor applications before general recursive definitions or induction tactics.}
Recursor skeletons substantially widen useful program and proof synthesis while
preserving the current termination story.  They reuse family descriptors,
typed branch goals, provider premises, and source obligations introduced by the
preceding architectural work.
\end{recommendation}

\section{A staged dependent and indexed fragment}
\label{sec:dependent}

\subsection{Why a staged approach is feasible}

Full dependent program synthesis is not a single feature.  It combines
higher-order unification, definitional equality, index refinement, motive
synthesis, proof search, coercions, and source elaboration.  Attempting all of
that in \Djex{} would destroy the present division of responsibility.

The typed-skeleton architecture permits several useful fragments in increasing
order of difficulty.

\subsection{Stage D0: exact dependent transport}

This stage largely exists already through alpha-aware opaque atoms.  Make it
explicit and more reliable:

\begin{itemize}
  \item hash-cons dependent source atoms by an exact frontend fingerprint;
  \item distinguish definitional equality, alpha equality, and mere pretty-text
    equality;
  \item preserve occurrence IDs and binder scopes;
  \item let structural rules transport, pair, inject, select, and compose these
    atoms without inspecting their internals; and
  \item report hidden dependent atoms in projection-loss diagnostics.
\end{itemize}

This improves higher-order plumbing over propositions such as
\(\forall n, P\,n\) without adding dependent elimination.

\subsection{Stage D1: dependent pairs and subtypes as certified structures}

\code{Sigma}, dependent records, and \code{Subtype} can be represented by
family descriptors whose later field types refer to earlier branch locals via
opaque source tokens.  The engine need not substitute those source terms
itself.  It needs typed pattern binders and a scoped obligation environment.

A constructor skeleton for \(\Sigma x:A, B\,x\) has two subgoals:

\[
  a : A, \qquad b : B\,a.
\]

The second expected type is a source handle instantiated by the frontend after
\(a\) is chosen.  In the initial implementation, require \(a\) to be an exact
local/global term known to the frontend, not an arbitrary generated expression.
This still enables useful packaging and projection terms.

Elimination exposes the first component and a second component whose exact type
is supplied by the branch descriptor.  Lean verifies the dependency.

\subsection{Stage D2: indexed constructor refinement}

For an indexed inductive, constructor matching refines indices.  A descriptor
can attach branch-local equations and specialized field types to each
constructor.  Search proceeds structurally under a \code{RefinementContext}:

\begin{lstlisting}[language=Haskell]
data RefinementContext source = RefinementContext
  { refinedIndices :: [(source, source)]
  , localEqualities:: [ObligationId]
  , sourceScope    :: SourceScopeId
  }
\end{lstlisting}

The neutral checker ensures that a branch cannot use a refinement token outside
its scope.  The frontend decides whether refinements are definitional, require
an equality proof, or need an explicit transport.  A failed source check drops
the candidate.

Good initial targets include:

\begin{itemize}
  \item equality itself, with reflexivity and one controlled transport rule;
  \item \code{Fin} and simple zero/successor index splits;
  \item vectors under one constructor match; and
  \item user GADTs whose branch result indices are syntactically constructor
    forms after elaboration.
\end{itemize}

\subsection{Stage D3: motive holes and recursors}

Only after D1 and D2 should the engine generate dependent recursor skeletons.
The motive is an explicit frontend obligation with a neutral shape and exact
source expected sort.  Search can use common templates:

\begin{itemize}
  \item constant motive;
  \item target-as-motive after abstracting the major premise;
  \item one index-variable abstraction; and
  \item a frontend-supplied finite motive candidate table.
\end{itemize}

Lean elaboration determines universe correctness and dependent conversion.  No
negative completeness claim should cover an omitted motive template.

\subsection{What should remain opaque}

Even after these stages, the following should remain source-owned unless a
specific checked rule is added:

\begin{itemize}
  \item arbitrary term-level computation inside indices;
  \item unrestricted higher-order unification;
  \item synthesis of arbitrary equality lemmas;
  \item quotient and heterogeneous equality reasoning;
  \item well-founded recursion and termination measures;
  \item and source-specific coercion search beyond finite obligations.
\end{itemize}

This boundary is not a failure.  It lets the system automate a much larger
class of complex types while preserving an understandable and testable search
core.

\section{Semantic environment retrieval at project and Mathlib scale}
\label{sec:retrieval}

\subsection{Separate retrieval from search}

The current live-provider inventory is intentionally bounded and goal-relevant,
with project ratings and a semantic cache.  This validates the basic source
integration.  Scaling it should not mean handing thousands of declarations to
\Exference{} and hoping its ranking compensates.

Introduce a distinct \term{retrieval layer}.  Its output is still a finite,
checked provider package, so the existing engine boundary and soundness model
remain intact.  Retrieval can evolve independently and can be tested against
known useful provider sets.

\subsection{Persistent signature index}

For each declaration, store a compact normalized signature record:

\begin{lstlisting}[language=Haskell]
data ProviderIndexEntry = ProviderIndexEntry
  { providerId          :: GlobalId
  , providerResultHeads :: [HeadFeature]
  , providerArgumentHeads:: [HeadFeature]
  , providerForallShape :: ForallShape
  , providerArrowDepth  :: Int
  , providerFamilyUses  :: [FamilyId]
  , providerConstraints :: [ConstraintFeature]
  , providerSortClass   :: SortClass
  , providerTypeSize    :: Int
  , providerNamespace   :: NamespaceFeature
  , providerModule      :: ModuleId
  , providerBaseRating  :: Rating
  , providerFingerprint :: TypeFingerprint
  }
\end{lstlisting}

A \code{HeadFeature} should distinguish exact family identity from opaque
nominal text.  Include partial higher-kinded heads and structural family IDs,
so a query for \code{Prod} can retrieve declarations indexed under the same
family descriptor as saturated products.

In \Leant{}, the index can live in the synthesis companion already associated
with snapshots, or in a project cache keyed by Lean toolchain, imported module
fingerprints, and relevant options.  Incremental updates should append entries
for session declarations and invalidate only affected namespace/result-head
buckets.

\subsection{Backward type-reachability expansion}

Start from the exact goal feature set.  Retrieve providers whose result can
match the goal under cheap head/kind/arity filters.  Their unsolved argument
features become the next frontier.  Continue for a small semantic radius:

\[
  G_0 = \mathsf{features}(\text{goal}), \qquad
  G_{i+1} = G_i \cup
    \bigcup_{p : \mathsf{result}(p)\leadsto G_i}
      \mathsf{arguments}(p).
\]

This is not full unification.  It is an admissible retrieval over-approximation.
The checked engine still validates and searches exact types.  Useful filters
include:

\begin{itemize}
  \item exact result head and structural family identity;
  \item result-arrow depth;
  \item proper kind and sort compatibility;
  \item quantified-result shape;
  \item required local head features already present in the query;
  \item constraint satisfiability from the bounded evidence table;
  \item namespace/module proximity;
  \item and user ratings.
\end{itemize}

At each radius, apply diversity quotas: do not let dozens of near-identical
lemmas from one namespace crowd out a composition bridge.  Keep at least a few
entries per distinct result fingerprint, argument-head profile, and module.

\subsection{Proof-graph retrieval}

Successful synthesized terms provide valuable but deterministic local data.
Record a graph edge from a provider's result feature to each argument feature,
weighted by successful use count, final candidate rank, and verification
success.  Project-local learning can reorder retrieval while preserving an
unchanged lexical fallback and explicit ratings.

This is not machine learning in the opaque-model sense.  It is a transparent
weighted dependency graph.  Its cache can be deleted without changing
soundness or available rules; only ranking changes.

\subsection{Inventory packages and staged widening}

Replace raw widths \(1,4,16,80\) with semantic packages:

\begin{enumerate}
  \item exact result-head providers;
  \item one-step bridge providers;
  \item constraint/evidence-complete providers;
  \item two-step backward-reachable providers;
  \item namespace/project neighborhood;
  \item and the full bounded retrieved set.
\end{enumerate}

The current geometric widths can remain as caps within each package during
migration.  The important change is that every stage has a semantic reason and
can be cached by a query feature fingerprint.

\begin{recommendation}{Recommendation 6: build a persistent semantic index and backward relevance expander before raising the live-provider cap.}
This is the path to useful Mathlib-scale synthesis.  A larger unstructured
inventory would mostly increase queue pressure, duplicate candidates, and Lean
verification cost; a semantic retrieval layer increases the probability that a
small inventory contains the necessary composition chain.
\end{recommendation}

\section{Cooperative Djinn--Exference search}
\label{sec:cooperation}

\subsection{Keep independent semantics, share typed artifacts}

A single hybrid engine would risk losing the clearest property of the current
system: \Djinn{} has a complete logical interpretation for an admitted
fragment; \Exference{} has heuristic ranked search and evidence resolution.
The better design is a portfolio over shared typed artifacts.

Three levels of cooperation are increasingly powerful.

\subsubsection{Level 1: shared preparation and semantic caches}

Both engines consume the same interned types, source certificates, family
descriptors, normalized provider inventory, and plan fingerprints.  They share
candidate fingerprints and verification outcomes.  Search frontiers remain
independent.

This is low risk and should accompany the typed candidate graph.

\subsubsection{Level 2: progress-aware portfolio scheduling}

Replace fixed batch quotas after the compatibility prefix with a deterministic
portfolio scheduler.  Give the next slice of work to the engine with the best
recent marginal progress, measured by:

\begin{itemize}
  \item new semantic candidates produced;
  \item new subgoal fingerprints solved;
  \item reduction in unsolved obligation count;
  \item source-verification acceptance rate;
  \item unique family/provider coverage;
  \item and frontier growth per unit of engine work.
\end{itemize}

Retain minimum service guarantees: each standalone twelve-group opportunity
must still be reachable, and a complete \Djinn{} refutation cannot be delayed
indefinitely by an expanding \Exference{} queue.  Deterministic counters and
stable tie-breaking preserve reproducibility.

\subsubsection{Level 3: proof skeletons and leaf filling}

The typed candidate graph enables a more interesting composition:

\begin{enumerate}
  \item \Djinn{} produces a structural derivation with one or more opaque
    atomic/provider/evidence leaves.
  \item Each leaf becomes a typed subquery for \Exference{} with the exact
    local context and obligations.
  \item \Exference{} fills leaves using live providers, class evidence, or
    recursive-family rules.
  \item The combined graph is checked by the neutral checker and the source
    frontend.
\end{enumerate}

Conversely, an \Exference{} candidate can contain structural holes which
\Djinn{} solves completely.  This is especially promising for environment-rich
higher-order plumbing: \Exference{} chooses the right library bridge, while
\Djinn{} fills the pure structural lambda surrounding it.

The leaf protocol must prevent cyclic self-justification.  A subquery records
its parent candidate fingerprint and forbids a provider/certificate dependency
cycle unless the source family descriptor explicitly authorizes a recursor.

\subsection{Refutations in a cooperative portfolio}

Only a complete engine projection may refute its exact subgoal.  A filled
candidate can use such refutations to prune a branch locally, but the overall
query is refuted only when:

\begin{itemize}
  \item the root projection is complete;
  \item every structural plan required by the capability fingerprint has been
    exhausted;
  \item no omitted provider or source obligation could add proof power; and
  \item every cooperative leaf is either solved or refuted in a complete
    subfragment.
\end{itemize}

In practice, most new cooperative modes will be positive-only at first.  Their
presence simply adds a \code{ProjectionLoss} to any otherwise complete
provider-free refutation.  This is the conservative rule already used by many
current extensions, generalized structurally.

\section{Concrete changes by repository}
\label{sec:repo-changes}

\subsection{Djex responsibilities}

The following additions belong naturally in \Djex{} because they are
frontend-neutral and should be shared by both engines.

\paragraph{New neutral modules.}

\begin{itemize}
  \item \code{Language.Haskell.Synthesis.TypedGenerated}: typed term graph,
    patterns, application witnesses, obligations, and fingerprints;
  \item \code{Language.Haskell.Synthesis.Certificate}: checked source
    certificates, ownership, provenance bounds, and fidelity witnesses;
  \item \code{Language.Haskell.Synthesis.Family}: family/constructor/eliminator
    descriptors and recursion metadata;
  \item \code{Language.Haskell.Synthesis.Capability}: projection losses,
    capability fingerprints, and structured negative evidence;
  \item \code{Language.Haskell.Synthesis.Plan}: occurrence IDs, normalized
    plan lattice, deterministic adaptive scheduler, and plan fingerprints;
  \item \code{Language.Haskell.Synthesis.Fingerprint}: interned type, context,
    subgoal, and candidate identities; and
  \item optionally \code{Language.Haskell.Djex.Portfolio}: shared preparation
    and progress-aware orchestration without merging engine internals.
\end{itemize}

\paragraph{Backend adapters.}

Each engine adapter should produce a typed candidate derivation.  This need not
rewrite the engines immediately.  The first version can instrument existing
search/checker steps:

\begin{itemize}
  \item \Djinn{} maps proof-term constructors and formula nodes to typed graph
    nodes while erasing synthetic instantiation premises through explicit
    substitution witnesses;
  \item \Exference{} maps holes, inventory uses, unification substitutions,
    tuple/constructor steps, and case rules to the same graph;
  \item the current \code{Expression} is generated from the typed graph by a
    compatibility projection; and
  \item a small neutral checker validates every candidate before it leaves
    \Djex{}.
\end{itemize}

\paragraph{Public query API.}

Add a versioned request extension rather than changing every existing field:

\begin{lstlisting}[language=Haskell]
data QuerySemanticPackage variable source = QuerySemanticPackage
  { semanticFamilies     :: [FamilyDescriptor ...]
  , semanticCertificates :: [SubstitutionCertificate ...]
  , semanticOccurrences  :: OccurrenceTable source
  , semanticCapabilities :: RequestedCapabilities
  , semanticLimits       :: SemanticLimits
  }
\end{lstlisting}

Legacy callers receive an empty package derived from existing declarations and
provider evidence.  This preserves the curated public facade and permits a
long migration.

\paragraph{Limits.}

Keep the current public constants as compatibility defaults.  Introduce a
\code{SemanticLimits} record whose fields are measured in unique plan views,
certificate nodes, typed graph nodes, obligation count, and fingerprint cache
size.  Reject malformed lazy/cyclic values at sealing time exactly as current
kind and assignment bounds do.

\subsection{Leant responsibilities}

\Leant{} should remain the authority for elaborated Lean semantics.

\paragraph{Source package producer.}

Extend the generated Lean metaprogram to emit:

\begin{itemize}
  \item stable occurrence IDs for every serialized goal/provider/family node;
  \item exact binder roles, visibility, and sort-domain tokens;
  \item family descriptors derived from Lean inductive and recursor metadata;
  \item active-instance and other substitution certificates;
  \item exact source AST handles or bounded structural rendering metadata;
  \item opaque term/index fingerprints with source equality classes; and
  \item an environment generation/fingerprint for cache invalidation.
\end{itemize}

The producer already computes much of this information in specialized forms.
The change is to emit one typed package rather than several parser-position
specific strings and maps.

\paragraph{Renderer becomes source completion.}

\code{Leant.Synth.Render} should gradually stop inferring types from untyped
syntax.  Its core operation becomes:

\begin{enumerate}
  \item traverse a typed candidate graph;
  \item look up exact source metadata by occurrence/certificate/family ID;
  \item materialize Lean syntax or an AST command;
  \item discharge obligations in a bounded source-completion phase;
  \item elaborate against the exact goal; and
  \item return either a verified term or a structured rejection.
\end{enumerate}

The existing fitting code remains as a compatibility path for legacy
\code{Expression} candidates until all engines emit enough typed data.
Regression tests should assert that new candidates use the typed path and do
not silently fall back.

\paragraph{Verification API.}

Extend the backend worker response with a small error taxonomy:

\begin{lstlisting}[language=Haskell]
data VerificationFailure
  = TypeMismatch OccurrenceId SourceTypeSummary SourceTypeSummary
  | UnsolvedInstance OccurrenceId SourceGoalSummary
  | UnsolvedMetavariable OccurrenceId
  | UniverseMismatch OccurrenceId
  | BinderVisibilityMismatch OccurrenceId
  | UnknownSourceIdentity SourceIdentity
  | IndexRefinementFailure OccurrenceId
  | TerminationFailure OccurrenceId
  | OtherElaborationFailure String
\end{lstlisting}

This information is diagnostic and heuristic.  The full Lean message can still
be retained for debugging.  The engine never trusts the category for
soundness.

\paragraph{Provider index.}

Add a persistent index builder and query feature extractor outside
\code{Leant.Synth.Engine}.  The engine continues to receive a finite checked
inventory.  Provider cache keys incorporate the index generation and semantic
retrieval stage rather than only a width.

\subsection{Shared testing responsibilities}

Every cross-repository capability should continue the current paired pattern:

\begin{enumerate}
  \item neutral \Djex{} unit/property tests;
  \item public-facade tests through both engines;
  \item generated Haskell accepted by GHC where applicable;
  \item \Leant{} pure serializer/adapter/renderer tests;
  \item malformed and over-bound controls;
  \item a Lean transcript covering \Djinn{}, \Exference{}, and combined mode;
  \item and a commit in \Leant{} pinning the exact reviewed \Djex{} revision.
\end{enumerate}

The typed architecture adds two valuable differential checks: the neutral
checker must accept every engine candidate, and the source verifier must accept
every displayed candidate.  A mismatch between those layers is a localized
boundary bug rather than an unexplained text-rendering failure.

\section{Worked target classes}
\label{sec:targets}

This section illustrates what the proposed architecture would make possible.
The examples are not promises that one generic algorithm solves all of them.
They identify the smallest new semantic mechanism needed beyond the current
system.

\subsection{Higher-rank provider composition with exact source metadata}

Consider a schematic provider family:

\begin{lstlisting}[language=Haskell]
select
  :: forall f g.
     Evidence f g
  => Box (f A B)
  -> (forall x. f x x -> g x)
  -> Result (g A)
\end{lstlisting}

and a source environment establishing one correlated choice
\(f := (,)\), \(g := \mathsf{Either}\;E\), with a contextual rank-N argument.
The recent structural higher-kinded work can preserve closely related cases,
but the complete path currently relies on parallel canonical/source
representations and renderer fitting.

Under the proposed architecture:

\begin{enumerate}
  \item the frontend emits one substitution certificate for \code{select};
  \item the certificate refers to the canonical pair/sum family descriptors and
    carries exact Lean rendering metadata;
  \item the typed application node records the selected substitution and the
    specialized type of each term argument;
  \item any class evidence is a source obligation associated with the same
    certificate; and
  \item the renderer performs no type reconstruction---it materializes the
    known specialization and asks Lean to discharge the evidence.
\end{enumerate}

This generalizes the latest \code{Prod}/\code{Sum} fixes to arbitrary families
whose descriptors explicitly declare an unsaturated structural identity.

\subsection{A fold-shaped list program}

A target such as

\[
  (\alpha \to \beta \to \beta) \to \beta \to
  \mathsf{List}\,\alpha \to \beta
\]

has the obvious \code{List.foldr}-shaped inhabitant.  Current live-provider
search may find the library function if it is retrieved.  A source-certified
recursor descriptor allows synthesis even without a provider indexed under the
exact target spelling:

\begin{enumerate}
  \item choose the \code{List} recursor because the major premise is a list and
    the target result is uniform;
  \item create nil- and cons-branch subgoals;
  \item expose the recursive result as a typed local in the cons branch;
  \item let \Djinn{} synthesize the structural application of the user's step
    function, or let \Exference{} choose a relevant provider; and
  \item render an exact \code{List.rec}, \code{List.foldr}, or source-preferred
    eliminator application supplied by the descriptor.
\end{enumerate}

No recursive binding is generated, and Lean's recursor is already termination
checked.

\subsection{A dependent pair projection}

For

\[
  (\Sigma x : A, B\,x) \to A,
\]

the engine needs only the first field type.  A dependent-family descriptor
makes the pair eliminable; the second field remains source-typed under the
first binder but can be unused.  The typed case node exposes
\(x:A\) and \(y:B\,x\); \Djinn{} returns \(x\).  No equality reasoning or
motive synthesis is required.

For

\[
  (\Sigma x : A, B\,x) \to \Sigma x : A, B\,x,
\]

the current opaque-transport route can already return the input.  The new
descriptor merely gives the engine an additional constructor/eliminator view
without sacrificing the exact opaque identity.

\subsection{One-layer indexed refinement}

For a vector-like family, consider a function which eliminates an impossible
constructor at length zero.  The source descriptor says that the nil
constructor returns index zero and cons returns successor.  In a branch where
the major premise has index zero:

\begin{itemize}
  \item the nil branch is admitted directly;
  \item the cons branch carries an equality obligation
    \(\mathsf{succ}(n)=0\);
  \item Lean can discharge it as impossible during elaboration or via a source
    no-confusion rule; and
  \item the candidate graph records why the branch is impossible rather than
    treating the whole indexed family as unrelated atoms.
\end{itemize}

This is a bounded constructor-refinement rule, not general dependent proof
search.

\subsection{Class-constrained higher-order plumbing}

Suppose the environment contains:

\begin{lstlisting}[language=Haskell]
mapMLike :: Monad m => (a -> m b) -> Container a -> m (Container b)
\end{lstlisting}

and the target supplies a function \code{a -> m b} and a container.  The engine
can structurally choose and apply the provider while leaving \code{Monad m} as
an evidence obligation.  In \Leant{}, the frontend asks type-class synthesis to
fill it in the exact context.  If no instance exists, the candidate is rejected
or retained as a visibly incomplete skeleton in a future diagnostic mode; it
is never displayed as a finished term.

This is more expressive than erasing dictionaries entirely, but safer and more
portable than importing source dictionary terms into the neutral engine.

\subsection{Library bridge plus structural completion}

A common complex synthesis shape is:

\[
\begin{aligned}
&\text{local higher-order plumbing}
  \quad\circ\quad \text{environment provider(s)}\\
&\hspace{7em}\circ\quad
  \text{local constructor/recursor plumbing}.
\end{aligned}
\]

A cooperative portfolio can solve this efficiently:

\begin{enumerate}
  \item semantic retrieval selects a provider whose result head reaches the
    goal;
  \item \Exference{} chooses that bridge and leaves typed argument holes;
  \item \Djinn{} fills purely structural holes completely;
  \item source evidence fills class/equality obligations; and
  \item semantic normalization removes administrative differences before Lean
    verification.
\end{enumerate}

This is the most plausible route to substantially more automation for complex
types.  It uses each subsystem where it is strongest instead of forcing one
engine to solve the entire problem monolithically.

\section{Soundness, completeness, and trust boundaries}
\label{sec:soundness}

\subsection{Positive soundness remains source-verified}

The fundamental positive-soundness theorem is operational:

\begin{quote}
A candidate is presented as a term of source type \(T\) only after the source
frontend has elaborated and checked that exact candidate against \(T\) in the
current sealed environment.
\end{quote}

The typed neutral checker is defense in depth and a debugging aid.  It catches
engine/boundary inconsistencies earlier, permits semantic normalization, and
makes source failures more meaningful.  It does not replace GHC or Lean.

Residual obligations preserve this theorem.  A candidate with unfilled
obligations is a skeleton, not a finished term.  The source completion layer
must fill every obligation and re-check the result before display.

\subsection{Certificate trust is bounded and checked}

A source certificate is not trusted merely because it came from a frontend.
The sealed query validates the structural facts on which search relies:

\begin{itemize}
  \item no duplicate or stale owner identity;
  \item exact telescope fingerprint and argument count;
  \item well-formed kinds and closed substitutions;
  \item bounded source payload size;
  \item no free source variables outside the declared scope;
  \item no conflicting family arities or class kind vectors;
  \item and required fidelity through datatype fields when nominal identity is
    used to justify a visible specialization.
\end{itemize}

Source-specific claims that \Djex{} cannot validate are obligations, not search
axioms.  For example, ``Lean can synthesize this instance'' is confirmed by
Lean during source completion; it is not used to prove a \Djinn{} refutation.

\subsection{Negative evidence must be monotone under capability loss}

Let \(P\) be the exact source problem and \(\pi_C(P)\) a logical projection
under capability set \(C\).  A \Djinn{} refutation is source-valid only if the
projection is complete for the stated fragment.  Adding positive-only
certificates or omitted source rules must not strengthen that claim.

A useful invariant is monotonicity:

\[
  C_1 \subseteq C_2
  \quad\Longrightarrow\quad
  \mathsf{loss}(C_2) \subseteq \mathsf{loss}(C_1).
\]

A refutation at \(C_1\) can be reused at \(C_2\) only when the newly added
capabilities introduce no new proof power relevant to the root problem, or when
those capabilities are themselves exhaustively searched.  Otherwise the old
refutation is retained only as an internal provider-free fallback, exactly as
\Leant{} already does.

\subsection{Refutation certificates}

A future \Djinn{} adapter could emit a compact certificate of saturated LJT
search rather than only a status.  The certificate need not be a proof term of
negation; it can be a replayable finite derivation showing that every applicable
rule branch closes unsuccessfully under the normalized formula and context.

Benefits include:

\begin{itemize}
  \item independent checking of negative results;
  \item stable caching keyed by formula and capability fingerprint;
  \item clearer diagnostics for which atom or omitted rule prevents a source
    refutation; and
  \item safer cooperative pruning of typed subgoals.
\end{itemize}

This is lower priority than typed positive candidates because the current
\Djinn{} boundary already has a strong completeness discipline.  It becomes
more valuable once capabilities and cooperative subqueries are more numerous.

\subsection{No semantic inference from source spelling}

The current projects already move away from raw text: validated names,
alpha-normalized closed types, exact family keys, and structural context
payloads replace guessed strings.  The target architecture should make the
rule absolute:

\begin{quote}
Source spelling is for presentation and source AST construction.  Search
identity, kinds, family arity, binder roles, and evidence ownership come only
from checked structural data.
\end{quote}

This is particularly important for Lean, where pretty printing can omit
universes, implicit arguments, qualifications, coercions, and binder
information.

\section{Implementation roadmap}
\label{sec:roadmap}

The roadmap below is dependency-ordered.  Effort labels are relative
architectural sizes, not calendar estimates.

\subsection{Milestone 0: observability and baseline corpus (S/M)}

Before changing the candidate edge, make current costs visible.

\begin{itemize}
  \item Add per-query counters for compiled rank-N plans, unique formula
    fingerprints, prepared-premise cache hits, engine states, raw derivations,
    rendered groups, semantic/exact duplicates, Lean variants tried, and
    verification failure classes.
  \item Record provider retrieval stage, inventory composition, and provider
    use in each verified candidate.
  \item Store machine-readable benchmark output beside the human transcripts.
  \item Freeze a representative corpus from the current pure tests and Lean
    goldens, including every known fail-closed sentinel.
\end{itemize}

This milestone produces an immediate result: it will reveal whether current
latency is dominated by plan compilation, \Exference{} queue growth, renderer
variant multiplication, duplicate candidate groups, or Lean round trips for
each class of query.

\subsection{Milestone 1: typed candidate spine and occurrence IDs (M/L)}

Implement the smallest typed slice that eliminates recent reconstruction work:

\begin{itemize}
  \item stable source occurrence IDs;
  \item typed local/global/application/type-application/tuple/let/case nodes;
  \item application witnesses containing the exact neutral function domain and
    result;
  \item provider certificate IDs on visible specializations;
  \item typed pattern fields for product and finite-family elimination;
  \item a neutral checker; and
  \item projection back to the existing \code{Expression}.
\end{itemize}

Migrate provider-result specialization and occurrence-local metadata to the
typed path first.  Keep the old renderer fitting as a fallback.  Success means
that the current provider structural-elimination, contextual-assignment,
implicit-forall, and higher-order-result regressions pass without using the
legacy inference path.

\subsection{Milestone 2: semantic fingerprinting and shared preparation (M)}

\begin{itemize}
  \item Intern neutral types, substitutions, contexts, and family IDs.
  \item Compute candidate fingerprints before rendering.
  \item Deduplicate each engine and combined mode by fingerprint.
  \item Share provider/certificate/family preparation across engines.
  \item Add dominance pruning over obligation/loss/cost vectors.
\end{itemize}

This milestone should reduce Lean verification attempts and make the existing
60/12/24 windows materially more valuable without changing search power.

\subsection{Milestone 3: adaptive occurrence-plan tail (L)}

\begin{itemize}
  \item Assign stable IDs to positive forall occurrences.
  \item Preserve the historical plan prefix exactly.
  \item Add normalized plan fingerprints and a best-first tail under a distinct
    compiled-view budget.
  \item Re-express the twelve-site six/six sentinel as an adaptive success.
  \item Add larger adversarial cases to demonstrate bounded behavior without
    adding fixed sextic code.
\end{itemize}

Structured verifier feedback can be added after a deterministic engine-only
heuristic proves the architecture.

\subsection{Milestone 4: generalized certificates and family descriptors (L)}

\begin{itemize}
  \item Lift provider assignments into a common certificate table while
    preserving the existing public constructors as adapters.
  \item Introduce family descriptors for current intrinsic and query-wide
    finite/recursive families.
  \item Move total arity, structural/nominal identity, constructor maps,
    phantom fidelity, and source rendering handles into descriptors.
  \item Reimplement \code{Prod}/\code{Sum}, tuples, \code{Either}, \code{List},
    and \code{Nat} through the common path.
\end{itemize}

No new user-visible family need be added until the current families pass
through the descriptor path.  This is an abstraction-validation milestone.

\subsection{Milestone 5: source obligations and recursor skeletons (L/XL)}

\begin{itemize}
  \item Generalize residual constraints to obligations.
  \item Add bounded evidence-table nodes and Lean source completion.
  \item Serialize non-dependent recursor descriptors for \code{Nat},
    \code{List}, and simple user inductives.
  \item Generate one-recursor branch skeletons and solve branches with existing
    engines.
  \item Mark omitted recursors as projection loss for negative evidence.
\end{itemize}

This milestone should produce the first qualitatively new class of synthesized
programs rather than only broader higher-rank transport.

\subsection{Milestone 6: dependent structures and indexed refinement (XL)}

Proceed through D0--D3 from Section~\ref{sec:dependent}.  Do not begin with
arbitrary dependent recursors.  The first acceptance criteria should be exact
dependent transport, \code{Sigma}/\code{Subtype} construction and projection,
and one-layer index-refining cases with source equality obligations.

\subsection{Milestone 7: semantic provider index and portfolio cooperation (L/XL)}

The retrieval index can be developed in parallel after stable family/type
fingerprints exist.  Progress-aware scheduling and leaf filling should wait for
typed candidate graphs and shared subgoal keys.

\begin{table}[H]
\centering
\small
\caption{Priority matrix.  Impact is on synthesis reach and architectural
leverage; risk is implementation/semantic risk.}
\begin{tabularx}{\textwidth}{L{4.7cm} C{1.35cm} C{1.35cm} C{1.55cm} Y}
\toprule
\textbf{Direction} & \textbf{Impact} & \textbf{Effort} & \textbf{Risk} &
\textbf{Dependency} \\
\midrule
Typed candidate graph + occurrence IDs & Very high & L & Medium & None \\
Semantic candidate fingerprinting & High & M & Low & Typed nodes \\
Adaptive rank-N plan tail & High & L & Medium & Occurrence IDs \\
General source certificates & Very high & L & Medium & Typed nodes \\
Family descriptor algebra & Very high & L & Medium & Certificates helpful \\
Evidence/equality obligations & High & L & Medium & Typed nodes \\
Recursor skeletons & Very high & L/XL & Medium & Families + obligations \\
Dependent/indexed refinement & Very high & XL & High & All preceding semantic layers \\
Persistent semantic provider index & High & L & Low--med. & Stable fingerprints \\
Cooperative leaf filling & High & XL & Med.-high & Typed DAG + shared cache \\
Refutation certificates & Medium & L & Medium & Capability model \\
\bottomrule
\end{tabularx}
\end{table}

\begin{recommendation}{Recommendation 7: make Milestone 1 the next architectural project and use one current provider-result regression as its vertical slice.}
A good first slice is a provider application whose exact first term argument
specializes a later rank-N domain, followed by structural elimination of the
provider result.  Today that path exercises mixed-spine recovery, substitution,
fitting, occurrence metadata, and explicit binder weaving.  Carrying the same
information in typed nodes would demonstrate the new interface against a
hard, already well-tested case.
\end{recommendation}

\section{Evaluation and regression strategy}
\label{sec:evaluation}

\subsection{Measure semantic work, not only wall time}

Interactive wall time matters, but it is a noisy aggregate of pure search,
subprocess scheduling, Lean elaboration, cache state, and machine load.  Each
benchmark should also report deterministic semantic counters:

\begin{itemize}
  \item source nodes serialized and opaque atoms introduced;
  \item exact family descriptors and certificates admitted/rejected;
  \item unique rank-N plans compiled versus raw plans proposed;
  \item unique formula and prepared-premise fingerprints;
  \item \Djinn{} sequents/rule branches and \Exference{} steps/queue prunes;
  \item shared subgoal cache hits;
  \item raw candidate derivations, semantic fingerprints, and dominated terms;
  \item source alternatives generated and verification attempts;
  \item verification outcomes by failure class;
  \item obligations created, solved, and rejected; and
  \item capability losses attached to the final outcome.
\end{itemize}

A change is an improvement if it reduces expensive semantic work or extends the
solved corpus under the same explicit budgets, even when local wall-clock noise
hides the gain.

\subsection{Benchmark dimensions}

Build a matrix rather than a flat list of showcase types.

\paragraph{Quantifier planning.}

Vary:

\begin{itemize}
  \item independent positive occurrence count from 1 through at least 20;
  \item required open/opaque balance, especially central antichains;
  \item nesting depth and shared ancestor paths;
  \item duplicated occurrences that compile to the same formula;
  \item and interaction with query-correlated or closed instantiations.
\end{itemize}

The current ten-, eleven-, and twelve-site sentinels should remain fixed anchor
points.  Add cases where many raw plans normalize to few formula views to test
semantic budgeting.

\paragraph{Provider certificates.}

Vary:

\begin{itemize}
  \item binder width, kinds, and vacuity;
  \item complete versus partial vectors;
  \item multiple vectors per provider;
  \item nested quantified/contextual arguments;
  \item canonical structural and nominal family heads;
  \item phantom, mixed faithful/phantom, recursive, and empty datatypes;
  \item occurrence-local rendering alternatives; and
  \item conflicting or stale source metadata.
\end{itemize}

Every positive should have a malformed or semantically unsafe neighbor which
fails closed without aborting unrelated evidence.

\paragraph{Candidate identity.}

Construct queries with many equivalent derivations:

\begin{itemize}
  \item repeated providers with alpha-equivalent types;
  \item eta-expanded and contracted forms;
  \item tuple construction followed by projection;
  \item source alternatives differing only in explicit binder spelling;
  \item equivalent synthetic premise paths; and
  \item \Djinn{}/\Exference{} candidates normalizing to the same graph.
\end{itemize}

Assert the number and order of semantic fingerprints separately from rendered
spellings.

\paragraph{Recursive and dependent stages.}

For recursor work, vary constructor count, branch arity, recursive fields,
independent SCCs, and nested recursor depth.  For dependent stages, vary exact
transport, dependent pair construction/projection, index-refining cases,
definitional versus propositional equality, and deliberately unsupported
motives.

\paragraph{Environment retrieval.}

Use synthetic environments containing controlled distractors as well as real
Lean projects.  Record recall of a known minimal provider set at each semantic
retrieval stage, inventory size, search work, and verified result rank.

\subsection{Property-based and metamorphic tests}

The current hand-crafted regressions are strong.  Typed interfaces enable
additional systematic properties.

\begin{description}
  \item[Alpha invariance.]
    Renaming source binders or private neutral identities does not change
    semantic candidate fingerprints or outcomes, modulo source display names.

  \item[Occurrence separation.]
    Duplicating a provider use creates distinct occurrence IDs and permits
    independent metadata choices without conflating substitutions.

  \item[Projection round trip.]
    Typed candidate \(\to\) legacy \code{Expression} \(\to\) typed checking
    preserves the root neutral type for the compatibility fragment.

  \item[Certificate permutation resistance.]
    Reordering certificate storage does not change semantic results; reordering
    the telescope inside a certificate is rejected or changes its fingerprint.

  \item[Bound productivity.]
    Infinite, cyclic, or excessively wide caller-built lists/kinds/graphs are
    rejected without forcing beyond the advertised bound.

  \item[Prefix preservation.]
    Enabling an adaptive tail, source obligations, or cooperative search leaves
    every historical candidate prefix byte-for-byte or fingerprint-for-fingerprint
    unchanged until the old frontier is exhausted.

  \item[Capability monotonicity.]
    Adding a positive-only capability can add candidates or weaken a negative
    verdict, but cannot create a stronger refutation.

  \item[Frontend independence.]
    Source rendering metadata changes spelling but not neutral typing or search
    identity unless it is attached to a semantically distinct certificate.
\end{description}

\subsection{Fuzzing the checked boundary}

The certificate/family/typed-graph parsers are security-like boundaries even in
a local tool: they consume untrusted caller-built lazy Haskell values and
machine-generated source metadata.  Add bounded generators for:

\begin{itemize}
  \item cyclic and shared graphs;
  \item duplicate IDs and dangling references;
  \item malformed kinds and arity conflicts;
  \item scope escape and binder capture;
  \item inconsistent field types;
  \item certificate owners from another session;
  \item branch refinements escaping their branch;
  \item and terms whose annotation disagrees with the neutral checker.
\end{itemize}

The expected result is deterministic rejection with a focused diagnostic, not
an exception, nontermination, or contamination of sibling evidence.

\subsection{Golden tests without freezing accidents}

Keep exact golden transcripts for source-visible behavior, especially candidate
ordering and explanatory notes.  For internal architectural changes, add a
parallel semantic golden format containing fingerprints, obligations, losses,
and normalized term graphs.  This avoids forcing incidental private names or
administrative lambdas when the user-visible output remains equivalent.

\section{Risks and mitigations}
\label{sec:risks}

\subsection{The typed graph could become a second source language}

\textbf{Risk.}  Adding binders, sorts, indices, equality, recursors, and source
metadata may grow the neutral IR until it resembles Lean's full expression
language.

\textbf{Mitigation.}  Keep source terms and indices opaque; add only distinctions
required by a generic search rule; make every source-specific operation an
obligation or certificate; and require at least two frontend uses, or a clear
engine-independent invariant, before moving a fact into the shared core.

\subsection{Two representations can drift}

\textbf{Risk.}  During migration, \code{Frag}, shared \code{Type}, typed
candidates, and legacy \code{Expression} may encode overlapping semantics.

\textbf{Mitigation.}  Define one-way checked projections and round-trip
properties.  Do not maintain independent hand-written interpretations in each
engine.  Mark compatibility adapters explicitly and migrate one vertical slice
at a time.

\subsection{Adaptive planning could destabilize candidate order}

\textbf{Risk.}  A best-first tail changes which large-plan candidate appears
first and may be sensitive to heuristic details.

\textbf{Mitigation.}  Preserve the complete historical prefix; use stable
source-order tie-breaking; version the tail policy; expose deterministic plan
fingerprints in diagnostics; and treat heuristic order beyond the old frontier
as experimental.

\subsection{Source feedback could accidentally affect soundness}

\textbf{Risk.}  Structured Lean failures used for refinement may be mistaken
for proof that a plan is impossible.

\textbf{Mitigation.}  Feedback changes only positive-search priority.  It may
never eliminate a completeness-required plan or strengthen negative evidence.
Capability and refutation code should not depend on verifier failure classes.

\subsection{Obligation search could explode}

\textbf{Risk.}  Allowing class, equality, coercion, and universe holes may turn
one term skeleton into a huge source completion space.

\textbf{Mitigation.}  Seal finite evidence tables; cap obligations and source
alternatives per candidate; rank obligation-free candidates first; deduplicate
by obligation fingerprint; and use one outer command deadline exactly as the
current staged searches do.

\subsection{Family descriptors may trust too much frontend data}

\textbf{Risk.}  A malformed descriptor could give an engine an unsound
constructor or eliminator rule.

\textbf{Mitigation.}  Validate all neutral structural invariants; treat
source-only facts as obligations; independently check every positive candidate;
and exclude descriptor-dependent refutations until the descriptor class has a
checked completeness proof.  Lean-derived descriptors can include a compact
source fingerprint and be regenerated after environment changes.

\subsection{Persistent retrieval can become nondeterministic}

\textbf{Risk.}  Learned weights or cache history may make results hard to
reproduce.

\textbf{Mitigation.}  Store transparent counts; version index formats; include
index generation in transcripts; provide a lexical/rating-only mode; and make
all tie-breaking stable.  Cache deletion may affect speed and ranking, never
soundness.

\section{Directions not recommended as the primary next step}
\label{sec:not-recommended}

\begin{caution}{Do not make the next major project merely ``raise six to seven.''}
The shared bound is well-factored, so a seven-binder widening may be a cheap and
useful incremental change.  It does not address the semantic recovery tax,
quantifier-plan lattice, dependent structure, or environment composition.  It
should be treated as maintenance, not the strategic roadmap.
\end{caution}

\begin{caution}{Do not add sextic, septic, and higher fixed frontiers indefinitely.}
Keep the quintic family as a deterministic prefix and validation oracle.  Put
new coverage in a normalized adaptive tail whose budget is unique compiled
views.  Otherwise each central antichain creates another bespoke scheduler and
documentation burden.
\end{caution}

\begin{caution}{Do not move raw Lean syntax into Djex as semantic authority.}
Exact source text is fragile under qualification, pretty-printing options,
implicit arguments, binder visibility, universes, and notation.  Carry stable
IDs and structural metadata; let \Leant{} build Lean syntax at the final edge.
\end{caution}

\begin{caution}{Do not rely on blind generate-and-reject as the main dependent-type strategy.}
Lean verification makes blind search safe, not effective.  Explicit typed
skeletons and obligations should constrain candidate generation before the
expensive elaborator call.
\end{caution}

\begin{caution}{Do not merge Djinn and Exference into one monolithic algorithm.}
Their independent completeness and ranking semantics are valuable.  Share
interned artifacts, candidate identity, and typed leaves; cooperate through a
portfolio protocol instead of erasing the distinction.
\end{caution}

\begin{caution}{Do not expose stronger non-inhabitation wording merely because the source checker rejected every generated variant.}
Rejected candidates are positive-search evidence only.  Negative evidence must
continue to come from a complete logical projection with an explicit capability
fingerprint.
\end{caution}

\section{Overall assessment}

The recent development of \Djex{} and \Leant{} is exceptionally coherent.  The
repositories are not duplicating work.  They form a two-stage research and
engineering system:

\begin{itemize}
  \item \Leant{} supplies difficult, source-realistic Lean cases and an
    independent semantic oracle;
  \item \Djex{} converts the reusable part of each case into a checked,
    frontend-neutral search capability;
  \item \Leant{} then proves that the abstraction was neither too weak to render
    nor too permissive to trust.
\end{itemize}

That method should continue.  The architectural unit of progress, however,
needs to change.  The last several capability loops increasingly preserve
semantic provenance through parallel maps and recover types after search.  The
next abstraction should make provenance, expected types, substitutions,
family identity, and residual source obligations first-class in the candidate
itself.

With that layer in place, the future directions reinforce one another:

\begin{enumerate}
  \item typed nodes enable semantic deduplication and shared subgoal caches;
  \item stable occurrence IDs enable adaptive quantifier planning and exact
    source metadata;
  \item generalized certificates unify provider, class, family, and equality
    evidence;
  \item family descriptors support canonical higher-kinded identity, recursive
    recursors, and indexed refinement;
  \item obligations let Lean solve source-specific evidence without infecting
    the neutral engine;
  \item semantic retrieval finds the small provider set needed for complex
    compositions; and
  \item a portfolio can let \Djinn{} solve structural leaves while
    \Exference{} selects environment bridges.
\end{enumerate}

The likely result is not a universal synthesizer for arbitrary Lean types.
It is something more useful and technically defensible: a substantially larger,
explicitly bounded class of complex higher-order, higher-rank, constrained,
recursive, and lightly dependent types for which the system can construct
verified terms automatically, explain what semantic resources it used, and say
precisely when a negative result is or is not justified.

\begin{recommendation}{Final recommendation}
Make a typed candidate graph with occurrence identity, substitution provenance,
and residual obligations the next shared architectural milestone.  Validate it
by migrating one of the current hardest provider-result/structural-elimination
paths.  Then use that foundation for semantic deduplication, an adaptive rank-N
plan tail, generalized source certificates, and family descriptors.  Those four
steps convert the successful Djex--Leant feedback loop from a sequence of
careful special cases into a platform for the next generation of synthesis
rules.
\end{recommendation}

\appendix

\section{Representative commit map}
\label{app:commits}

The following list is not exhaustive.  It records the commits most directly
supporting the architectural inferences in this report.

\subsection{Djex}

\begin{longtable}{L{2.3cm} L{4.7cm} L{7.2cm}}
\toprule
\textbf{Commit} & \textbf{Theme} & \textbf{Architectural significance} \\
\midrule
\endfirsthead
\toprule
\textbf{Commit} & \textbf{Theme} & \textbf{Architectural significance} \\
\midrule
\endhead
\bottomrule
\endfoot
\djcommit{d728719ff2515a8a464398aaab3e789420cc86a6}{d728719f} & Quintic rank-N frontiers & Fixed-degree plan family reaches complete coverage through eleven independent sites under a finite view cap. \\
\djcommit{227bdc52e66e0b07606b84e89a7c3516493a718b}{227bdc52} & Closed local instantiation & Positive-only query-local scheme specialization at closed query subtrees. \\
\djcommit{077f767fd75814bbcd74bed795e10455935e397f}{077f767f} & Query-correlated instantiation & Chooses supplied quantified types only when the complete specialized body occurs in the query. \\
\djcommit{db25170b2aed13d2bb1249c11229cc28f809aa95}{db25170b} & Scalar aggregate elimination & Exposes checked tuple structure from globals and provider evidence to \Exference{}. \\
\djcommit{e10f7f018facef4445fc774118b39d716b36c02e}{e10f7f01} & Mixed application spine & One shared constructor for term and visible-type arguments, preserving order and wide-spine strictness. \\
\djcommit{03aa3820284f9117011910b7eb931568ef4255cd}{03aa3820} & Contextual exact assignments & Closed contextual rank-N vectors with occurrence-sensitive fidelity marking. \\
\djcommit{103fdd391aaa350c09f07b682d6fd70a67a44989}{103fdd39} & Sibling-field fidelity & Prevents one faithful field from masking a selected argument erased in another field. \\
\djcommit{c5afc68a1b49adfb66f25d67c49402487cf9697d}{c5afc68a} & Kinded shared classes & Preserves authoritative class-parameter kinds at the neutral environment boundary. \\
\djcommit{66955da710951fd525f7ad874d133132ffcd2cee}{66955da7} & Six-binder widening & One public bound advances hypothesis, loaded-scheme, and provider routes together. \\
\djcommit{5a550eb409c37c00256a41b14bad5a0d3792fa68}{5a550eb4} & Structural higher-kinded assignment & Connects bare pair-constructor evidence to the same identity as saturated tuple structure. \\
\djcommit{f50b5eac86726c7530d95e48c27101b7c0fa25c9}{f50b5eac} & Reviewed head & Merges canonical structural higher-kinded assignment support. \\
\end{longtable}

\subsection{Leant}

\begin{longtable}{L{2.3cm} L{4.7cm} L{7.2cm}}
\toprule
\textbf{Commit} & \textbf{Theme} & \textbf{Architectural significance} \\
\midrule
\endfirsthead
\toprule
\textbf{Commit} & \textbf{Theme} & \textbf{Architectural significance} \\
\midrule
\endhead
\bottomrule
\endfoot
\lecommit{e09bded3de507c2b28bb243421cdfe07371aaa1d}{e09bded3} & Explicit provider heads & Reconstructs explicit forall placeholders for discovered provider applications. \\
\lecommit{15fb0962ca7cf99c6df6a8c0c0776d79febe7ac9}{15fb0962} & Nested fitting traversal & Reaches provider applications beneath lambdas, tuples, unknown applications, and eliminations. \\
\lecommit{690277ca282b8c196b6a8c34777e8689539cd582}{690277ca} & Rank-N fields through elimination & Transfers exact provider result fragments into case/let-bound fields. \\
\lecommit{892117dcc9756f3ede4b8500665892781e679fb7}{892117dc} & Result specialization from terms & Infers quantified fragment substitutions from exact term arguments. \\
\lecommit{77ee65b36a0544b3ce5cc37f44a6a3eb9d76caff}{77ee65b3} & Full-spine specialization & Propagates learned substitutions to later higher-rank arguments. \\
\lecommit{888e2b402529089e606a55c8374d9bbf1a7568d4}{888e2b40} & Exact forall rendering metadata & Retains visibility and \code{Prop}/\code{Type}/\code{Sort} domain tags beside canonical fragments. \\
\lecommit{fe4b2932d5ff356ac5792e26f698444d1b99fe81}{fe4b2932} & Occurrence-coupled metadata & Ensures the metadata used to fit one provider occurrence is the metadata later rendered there. \\
\lecommit{2979dbe2fe25915cdaf497dca96b934a85c6b4c1}{2979dbe2} & Structured contextual assignments & Carries exact nominal Lean class contexts through private engine classes and source-faithful rendering. \\
\lecommit{1668c956b2488bfeb4af2659314a0317a8e97fb5}{1668c956} & Higher-kinded context arguments & Preserves residual kind arity and total nominal-family identity. \\
\lecommit{d04c1a8fb9971cd7eec07fc873016627ebb1e928}{d04c1a8f} & Six-binder integration & Derives every bridge width from the shared \Djex{} authority and verifies all-engine Lean cases. \\
\lecommit{0dbe08a7819f44c872bf091685f6b30f170f0133}{0dbe08a7} & Structural higher-kinded synthesis & Separates canonical nominal evidence from legacy fragment payloads and maps \code{Prod}/\code{Sum} structurally. \\
\lecommit{0160237e4bdf6e342a55f5a86ce5d5292728024f}{0160237e} & Reviewed head & Merges structural higher-kinded synthesis and pins \Djex{} at \code{f50b5eac}. \\
\end{longtable}

\section{Suggested initial API migration}
\label{app:api}

A deliberately incremental API sequence is:

\begin{enumerate}
  \item Add \code{OccurrenceId}, \code{CertificateId}, and \code{FamilyId} as
    validated, collision-safe neutral identifiers.
  \item Add \code{TypedExpression} as an annotation wrapper around current
    \code{Expression} nodes rather than replacing the syntax immediately.
  \item Let engine details optionally carry a \code{TypedCandidate} and keep
    \code{candidateOutput} unchanged.
  \item Add a neutral checker and semantic fingerprint for the typed path.
  \item Extend \Leant{}'s provider maps to be keyed directly by occurrence and
    certificate IDs.
  \item Migrate one application/elimination vertical slice.
  \item Once all public engines emit typed nodes for a constructor, make the
    legacy renderer inference for that constructor diagnostic-only.
  \item Generalize residual constraints to obligations under a new API version.
  \item Introduce family descriptors and adaptive plans after typed occurrence
    IDs are stable.
\end{enumerate}

A compatibility envelope might be:

\begin{lstlisting}[language=Haskell]
data CandidateDetails backend ty source = CandidateDetails
  { backendDetails :: backend
  , typedCandidate :: Maybe (TypedCandidate ty LocalId source)
  , capabilityInfo :: CapabilityInfo source
  }
\end{lstlisting}

Existing callers ignore \code{typedCandidate}; \Leant{} opts in immediately;
new neutral tooling can inspect it without knowing \Djinn{} or \Exference{}
internals.

\section{Seed benchmark types}
\label{app:benchmarks}

The following schematic families are suitable for a future benchmark corpus.
They should be instantiated in both Haskell-like neutral tests and Lean goldens
where the source language permits.

\begin{description}
  \item[Central rank-N antichains.]
    Twelve or more independent positive quantified sites where the witness
    requires a balanced open/opaque choice not present in the fixed quintic
    prefix.

  \item[Correlated higher-kinded assignments.]
    Providers with two to six binders assigned to a mixture of \code{Prod},
    partially applied \code{Sum}, user constructors, and closed quantified
    types, including exact contextual class arguments.

  \item[Metadata-sensitive repeated uses.]
    The same provider used twice with alpha-equivalent canonical types but
    different explicit/implicit source binder spellings.

  \item[Higher-order result specialization.]
    An early exact term argument determines a later rank-N function domain and
    the provider returns a product whose polymorphic field is eliminated.

  \item[Phantom-fidelity adversaries.]
    Nested aliases and mixed faithful/phantom sibling fields which would permit
    an invalid visible specialization if any occurrence marker were lost.

  \item[Recursor skeletons.]
    \code{Nat} iteration, \code{List} fold, a binary tree catamorphism, and a
    two-constructor user family with one recursive field.

  \item[Dependent structures.]
    \code{Sigma}/\code{Subtype} projection and reconstruction, followed by
    one-layer indexed cases whose impossible branch requires an equality
    obligation.

  \item[Environment bridges.]
    Synthetic libraries where the solution requires exactly one, two, or three
    providers separated by distractors with the same result head.

  \item[Portfolio leaf filling.]
    A structural \Djinn{} skeleton with one provider leaf; an \Exference{}
    provider choice with two pure structural argument holes; and a mixed case
    with a source evidence obligation.
\end{description}

\section{Primary repository sources}
\label{app:sources}

The analysis used the following current repository materials in addition to the
commit history linked above:

\begin{itemize}
  \item \href{https://github.com/VladimirReshetnikov/Djex/blob/main/README.md}{Djex README}
    and its linked architecture, API, REPL, and dated implementation reports;
  \item \href{https://github.com/VladimirReshetnikov/Djex/blob/main/synthesis/src/Language/Haskell/Synthesis/Type.hs}{shared type representation};
  \item \href{https://github.com/VladimirReshetnikov/Djex/blob/main/synthesis/src/Language/Haskell/Synthesis/Generated.hs}{shared generated syntax};
  \item \href{https://github.com/VladimirReshetnikov/Djex/blob/main/synthesis/src/Language/Haskell/Synthesis/Candidate.hs}{candidate envelope};
  \item \href{https://github.com/VladimirReshetnikov/Djex/blob/main/synthesis/src/Language/Haskell/Synthesis/Query.hs}{query and provider-evidence bounds};
  \item \href{https://github.com/VladimirReshetnikov/Leant/blob/main/README.md}{Leant README};
  \item \href{https://github.com/VladimirReshetnikov/Leant/blob/main/docs/SYNTHESIS_PROPOSAL.md}{Leant synthesis proposal and implementation status};
  \item \href{https://github.com/VladimirReshetnikov/Leant/blob/main/src/Leant/Synth/Fragment.hs}{Lean fragment serializer/parser};
  \item \href{https://github.com/VladimirReshetnikov/Leant/blob/main/src/Leant/Synth/Engine.hs}{engine boundary and scheduling}; and
  \item \href{https://github.com/VladimirReshetnikov/Leant/blob/main/src/Leant/Synth/Render.hs}{source fitting, rendering, and candidate specialization}.
\end{itemize}

\vfill
\begin{center}
\small\sffamily
End of report.
\end{center}

\end{document}
